\subsection{卷积与线性表示}

命题 11. --- 设 G 为局部紧群，E 为拟完备局部凸空间，U 为 G 在 E 中的连续表示。

\begin{enumerate}
\def\labelenumi{(\roman{enumi})}
\item
  若 \(\lambda \in \mathcal{C}'(\mathbf{G})\)，\(\mu \in \mathcal{C}'(\mathbf{G})\)，则 \(U(\lambda * \mu) = U(\lambda)U(\mu)\)。
\item
  假设 \(\mathbf{E}\) 是 Banach 空间，并置 \(\rho(s) = \|U(s)\|\)（\(s\in\mathbf{G}\)）。若 \(\lambda\in\mathcal{M}^{\rho}(\mathbf{G})\)，\(\mu\in\mathcal{M}^{\rho}(\mathbf{G})\)，则 \(U(\lambda*\mu)=U(\lambda)U(\mu)\)。
\end{enumerate}

设 \(\lambda, \mu\) 属于 \(\mathcal{C}'(\mathrm{G})\)。对任意 \(x \in \mathrm{E}\)，应用第六章 §1, No.~1 的命题 1 和命题 4，即得

\[
\begin{array}{r l} & U (\lambda * \mu) x = \int_ {\mathbf {G}} U (s) x d (\lambda * \mu) (s) \\ & \quad = \int_ {\mathbf {G} \times \mathbf {G}} U (s t) x d \lambda (s) d \mu (t) = \int_ {\mathbf {G} \times \mathbf {G}} U (s) U (t) x d \lambda (s) d \mu (t) \\ & \quad = \int_ {\mathbf {G}} U (\lambda) U (t) x d \mu (t) = U (\lambda) \int_ {\mathbf {G}} U (t) x d \mu (t) = U (\lambda) U (\mu) x, \end{array}
\]

由此得 (i)。类似的论证可应用于情形 (ii)。

仍设 G 为局部紧群，并假设 G 连续地左作用于局部紧空间 X。这就定义了（§2, No.~4）G 的一个连续线性表示 \(\pmb{\gamma}\)，它作用在 \(\mathcal{M}(\mathrm{X})\) 中（后者赋予在 \(\mathcal{K}(\mathrm{X})\) 中的紧收敛拓扑）。

命题 12. --- 若 \(\lambda \in \mathcal{C}'(\mathbf{G})\) 且 \(\mu \in \mathcal{M}(\mathbf{X})\)，则

\[
\gamma (\lambda) \mu = \lambda * \mu .
\]

由 §1, No.~5 的命题 7，

\[
\lambda * \mu = \int_ {G} (\varepsilon_ {s} * \mu) d \lambda (s).
\]

而 \(\varepsilon_{s}*\mu = \pmb {\gamma}(s)\mu\)（No.~2 公式 (5)），且

\[
\int_ {\mathsf {G}} \left(\pmb {\gamma} (s) \mu\right) d \lambda (s) = \pmb {\gamma} (\lambda) \mu
\]

这是按 \(\pmb{\gamma}(\lambda)\) 的定义得到的。

推论. --- 映射 \((\lambda,\mu)\mapsto\lambda*\mu\)（从 \(\mathcal{C}^{\prime}(\mathrm{G})\times\mathcal{M}(\mathrm{X})\) 到 \(\mathcal{M}(\mathrm{X})\) 中）相对于 \(\mathcal{C}^{\prime}(\mathrm{G})\) 的等度连续子集与 \(\mathcal{M}(\mathrm{X})\) 的紧子集是亚连续的（\(\mathcal{C}^{\prime}(\mathrm{G})\) 与 \(\mathcal{M}(\mathrm{X})\) 分别赋予在 \(\mathcal{C}(\mathrm{G})\) 与 \(\mathcal{H}(\mathrm{X})\) 中的紧收敛拓扑）。

因为 \(\mathcal{M}(\mathrm{X})\) 赋予在 \(\mathcal{K}(\mathrm{X})\) 中的紧收敛拓扑后是拟完备的，所以映射 \((\lambda, \mu) \mapsto \pmb{\gamma}(\lambda)\mu\)（从 \(\mathcal{C}'(\mathrm{G}) \times \mathcal{M}(\mathrm{X})\) 到 \(\mathcal{M}(\mathrm{X})\) 中）相对于 \(\mathcal{C}'(\mathrm{G})\) 的等度连续子集与 \(\mathcal{M}(\mathrm{X})\) 的紧子集是亚连续的（§2, No.~6）。于是只需应用命题 12。

备注. --- 1) 设 \(\lambda_0 \in \mathcal{C}'(\mathbf{G})\)。映射 \(\mu \mapsto \lambda_0 * \mu\)（从 \(\mathcal{M}(\mathbf{X})\) 到 \(\mathcal{M}(\mathbf{X})\) 中）是淡连续的。因为，设 \(f \in \mathcal{K}(\mathbf{X})\)。我们有

\[
\langle \lambda_ {0} * \mu , f \rangle = \int f (s x) d \lambda_ {0} (s) d \mu (x) = \langle \mu , g \rangle ,
\]

其中 \(g(x) = \int f(sx) d\lambda_0(s)\)。而 \(g\) 是连续的（第七章 §1, No.~1 引理 1）。另一方面，设 \(S\) 为 \(\lambda_0\) 的支集，\(K\) 为 \(f\) 的支集。条件 \(sx \in K\) 与 \(s \in S\) 蕴含 \(x \in S^{-1}K\)；因此 \(g\) 的支集包含于 \(S^{-1}K\)，从而 \(g \in \mathcal{K}(X)\)。于是 \(\langle \lambda_0 * \mu, f \rangle = \langle \mu, g \rangle\) 是 \(\mu\) 的淡连续函数，这就证明了我们的断言。

\begin{enumerate}
\def\labelenumi{\arabic{enumi})}
\setcounter{enumi}{1}
\tightlist
\item
  设 \(\mu_0 \in \mathcal{M}(\mathrm{X})\)。映射 \(\lambda \mapsto \lambda * \mu_0\)（从 \(\mathcal{C}'(\mathrm{G})\) 到 \(\mathcal{M}(\mathrm{X})\) 中）对于拓扑 \(\sigma(\mathcal{C}'(\mathrm{G}), \mathcal{C}(\mathrm{G}))\) 与 \(\sigma(\mathcal{M}(\mathrm{X}), \mathcal{K}(\mathrm{X}))\) 是连续的。因为，设 \(f \in \mathcal{K}(\mathrm{X})\)。置 \(h(s) = \int f(sx) d\mu_0(x)\)，则 \(\langle f, \lambda * \mu_0 \rangle = \langle h, \lambda \rangle\)，且 \(h \in \mathcal{C}(\mathrm{G})\)（第七章 §1, No.~1 引理 1）。
\end{enumerate}

命题 13. --- 映射 \((s,\mu)\mapsto\boldsymbol{\gamma}(s)\mu\)（从 \(\mathbf{G}\times\mathcal{M}_{+}(\mathbf{X})\) 到 \(\mathcal{M}_{+}(\mathbf{X})\) 中）当 X 上的正测度所成的集合 \(\mathcal{M}_{+}(\mathbf{X})\) 赋予淡拓扑时是连续的。

由于 \(\pmb{\gamma}(s)\mu = \pmb{\gamma}(ss_0^{-1})\pmb{\gamma}(s_0)\mu\)，根据备注 1，只需证明所考虑的映射在形如 \((e,\mu_0)\)（其中 \(\mu_0\in \mathcal{M}_+(\mathrm{X})\)）的点处的连续性。给定函数 \(f\in \mathcal{K}(\mathrm{X})\) 与数 \(\varepsilon >0\)，问题在于证明：存在 \(e\) 在 G 中的邻域 U 与 \(\mu_0\) 在 \(\mathcal{M}_+(\mathrm{X})\) 中的邻域 W，使得关系 \(s\in \mathrm{U}\)、\(\mu \in \mathrm{W}\) 蕴含

\[
\left| \int f (s x) d \mu (x) - \int f (x) d \mu_ {0} (x) \right| \leqslant \varepsilon .\tag{6}
\]

设 \(\mathbf{V}\) 为 \(\mathbf{K}\)（即 \(f\) 的支集）在 \(\mathbf{X}\) 中的一个紧邻域，并设 \(\varphi \in \mathcal{K}_{+}(\mathbf{X})\) 满足 \(\varphi(x) = 1\)（在 \(\mathbf{V}\) 上）；存在 \(\mathrm{W}_0\)（\(\mu_0\) 在 \(\mathcal{M}_{+}(\mathbf{X})\) 中的一个邻域），使得 \(a = \sup_{\mu \in \mathrm{W}_0} \mu(\mathrm{V})\) 为有限：只需对 \(\mathrm{W}_0\) 取 \(\mu \in \mathcal{M}_{+}(\mathbf{X})\) 中满足 \(|\langle \varphi, \mu - \mu_0 \rangle| \leqslant 1\) 者所成的集合即可。另一方面，由于映射 \((s, x) \mapsto sx\) 连续，存在 \(\mathrm{U}_0\)（\(e\) 在 \(\mathbf{G}\) 中的一个紧邻域），使得 \(s\mathrm{K} \subset \mathrm{V}\) 对一切 \(s \in \mathrm{U}_0\) 成立；于是函数 \((s, x) \mapsto f(sx)\) 在 \(\mathrm{U}_0 \times \mathrm{V}\) 中一致连续，故存在邻域 \(\mathrm{U} \subset \mathrm{U}_0\)（\(e\) 的一个邻域），使得 \(|f(sx) - f(x)| \leqslant \varepsilon / 2a\) 对一切 \(s \in \mathrm{U}\) 与 \(x \in \mathrm{V}\) 成立。因此，对 \(s \in \mathrm{U}\) 与 \(\mu \in \mathrm{W}_0\)，我们有

\[
\left| \int f (s x) d \mu (x) - \int f (x) d \mu (x) \right| \leqslant \varepsilon / 2;
\]

若 \(\mathbf{W} \subset \mathbf{W}_0\) 是 \(\mu_0\) 在 \(\mathcal{M}_{+}(\mathbf{X})\) 中的由测度 \(\mu \in \mathbf{W}_0\)（满足 \(|\int f(x) d\mu(x) - \int f(x) d\mu_0(x)| \leqslant \varepsilon/2\)）所组成的邻域，则 \(\mathbf{U}\) 与 \(\mathbf{W}\) 满足要求。

\section{测度与函数的卷积}

\subsection{一个测度与一个函数的卷积}

设 X 为局部紧空间，局部紧群 G 连续地左作用于其上。设 \(\beta\) 为 X 上的正测度，在 G 下拟不变。设 \(\chi\) 为 \textgreater0 且定义在 \(G \times X\) 上的函数，它对 \(\mathbf{G} \times \mathbf{X}\) 上的每个测度都可测，且对每个 \(s \in \mathbf{G}\)，\(\chi(s^{-1}, \cdot)\) 是 \(\pmb{\gamma}(s)\beta\) 关于 \(\beta\) 的一个密度：

\[
\pmb {\gamma} (s) \beta = \chi (s ^ {- 1}, \cdot) \cdot \beta ,\tag{1}
\]

按照第七章 §1, No.~1 的约定，上式可写成：

\[
d \beta (s x) = \chi (s, x) d \beta (x).\tag{1'}
\]

这些数据在 No.~1、No.~2、No.~3 中保持固定（No.~2 的备注 2 除外）。

回顾（§2, No.~5）：若 \(\chi\) 连续且 \(\beta\) 的支集为 \(X\)，则 \(\chi\) 是一个乘子。

设 \(f\) 是局部 \(\beta\)-可积的复值函数，定义在 \(X\) 上，并设 \(\mu\) 是 \(G\) 上的测度。对每个 \(s \in G\)，测度 \(\boldsymbol{\gamma}(s)(f \cdot \beta)\) 以 \(\beta\) 为基，因为 \(\beta\) 是拟不变的。因此，若 \(\mu\) 与 \(f \cdot \beta\) 可作卷积，则 \(\mu * (f \cdot \beta)\) 以 \(\beta\) 为基（§3, No.~2 命题 10）。

定义 1. --- 若 \(\mu\) 与 \(f \cdot \beta\) 可作卷积，则称 \(\mu\) 与 \(f\) 关于 \(\beta\) 可作卷积。\(\mu * (f \cdot \beta)\) 关于 \(\beta\) 的每个密度都称为 \(\mu\) 与 \(f\) 关于 \(\beta\) 的一个卷积，记作 \(\mu * ^{\beta} f\)。

当不致混淆时省略 \(\beta\)。G 上多个测度与 X 上一个函数的卷积以类似方式定义。

\(\mu\) 与 f 的各个卷积在局部 \(\beta\)-几乎处处相等。若 \(\beta\) 的支集为 X，且存在 \(\mu\) 与 f 的一个连续卷积，则它是唯一确定的；这时称它为 \(\mu\) 与 f 关于 \(\beta\) 的卷积。

设 \(s \in G\)，并设 \(f\) 是局部 \(\beta\)-可积的复值函数，定义在 \(X\) 上。则 \(\varepsilon_s\) 与 \(f\) 可作卷积，且

\[
\varepsilon_ {s} * (f \cdot \beta) = \boldsymbol {\gamma} (s) (f \cdot \beta) = (\boldsymbol {\gamma} (s) f) \cdot (\boldsymbol {\gamma} (s) \beta) = (\boldsymbol {\gamma} (s) f) \cdot \chi (s ^ {- 1}, \cdot) \cdot \beta ,
\]

因此

\[
\left(\varepsilon_ {s} * f\right) (x) = \chi \left(s ^ {- 1}, x\right) f \left(s ^ {- 1} x\right) = \left(\gamma_ {\chi} (s) f\right) (x)\tag{2}
\]

在局部 \(\beta\)-几乎处处成立。

引理 1. --- 设 \(\mu\) 是 \(G\) 上的测度。则 \(\chi\) 是局部 \((\mu \otimes \beta)\)-可积的，且 \(\mu \otimes \beta\) 在同胚 \((s, x) \mapsto (s, s^{-1}x)\)（从 \(G \times X\) 到 \(G \times X\)）下的像是 \(\chi \cdot (\mu \otimes \beta)\)。

可以假设 \(\mu \geqslant 0\)。设 \(\mathbf{F} \in \mathcal{K}_{+}(\mathbf{G} \times \mathbf{X})\)。则

\[
\begin{array}{l} \iint \mathrm{F} (s, s ^ {- 1} x) d \mu (s) d \beta (x) = \int d \mu (s) \int \mathrm{F} (s, s ^ {- 1} x) d \beta (x) \\ = \int d \mu (s) \int \mathrm{F} (s, x) d (\boldsymbol {\gamma} (s ^ {- 1}) \beta) (x) = \int d \mu (s) \int \mathrm{F} (s, x) \chi (s, x) d \beta (x). \end{array}
\]

而函数 \((s,x)\mapsto\mathrm{F}(s,x)\chi(s,x)\) 具有紧支集且是 \((\mu\otimes\beta)\)-可测的。由第五章 §8, No.~3 命题 7，前面的等式证明该函数是 \((\mu\otimes\beta)\)-可积的，且

\[
\iint \mathrm{F} (s, s ^ {- 1} x) d \mu (s) d \beta (x) = \iint \mathrm{F} (s, x) \chi (s, x) d \mu (s) d \beta (x).
\]

这同时证明了引理 1 的两个断言。

命题 1. --- 设 \(\mu\) 是 G 上的测度，\(f\) 是 X 上局部 \(\beta\)-可积的复值函数。假设函数 \(s \mapsto f(s^{-1}x)\chi(s^{-1}, x)\) 是本质 \(\mu\)-可积的，但一个局部 \(\beta\)-可忽略的 \(x\) 值集合除外，且函数 \(x \mapsto \int |f(s^{-1}x)| \chi(s^{-1}, x) d|\mu|(s)\)（它对 \(\beta\) 局部几乎处处有定义）是局部 \(\beta\)-可积的。则 \(\mu\) 与 \(f\) 可作卷积。

可以假设 \(f \geqslant 0\) 且 \(\mu \geqslant 0\)。设 \(h \in \mathcal{H}_{+}(X)\)。我们要证明函数 \((s, x) \mapsto h(sx)\) 对于 \(\mu \otimes (f \cdot \beta) = (1 \otimes f) \cdot (\mu \otimes \beta)\) 是本质可积的（第五章 §8, No.~5 命题 10），亦即 \(\iint^{\bullet} h(sx)f(x)d\mu(s)d\beta(x) < +\infty\)（第五章 §5, No.~3 命题 3）；显然只需证明：存在 \(a > 0\)，使得对每个紧子集 \(K\)（\(G\) 中的），

\[
\iint^ {\bullet} h (s x) f (x) \varphi_ {K} (s) d \mu (s) d \beta (x) \leqslant a.
\]

由引理 1，

\[
\begin{array}{r l} \iint^ {\bullet} h (s x) f (x) \varphi_ {K} (s) d \mu (s) d \beta (x) \\ & = \iint^ {*} h (x) f (s ^ {- 1} x) \varphi_ {K} (s) \chi (s ^ {- 1}, x) d \mu (s) d \beta (x). \end{array}
\]

而函数 \((s,x)\mapsto h(x)f(s^{-1}x)\varphi_{\mathrm{K}}(s)\chi(s^{-1},x)\) 是 \((\mu\otimes\beta)\)-可测的（引理 1）且具有紧支集。于是前一表达式等于（第五章 §8, No.~3 命题 7）

\[
\begin{array}{r l} & {\int^ {*} h (x) d \beta (x) \int^ {*} f (s ^ {- 1} x) \varphi_ {\mathrm{K}} (s) \chi (s ^ {- 1}, x) d \mu (s)} \\ & {\qquad \leqslant (\sup h) \int_ {\mathrm{S}} ^ {*} d \beta (x) \int^ {\bullet} f (s ^ {- 1} x) \chi (s ^ {- 1}, x) d \mu (s),} \end{array}
\]

其中 S 表示 h 的支集。命题得证。

命题 2. --- 设 \(\mu\) 是 G 上的测度，\(f\) 是 X 上局部 \(\beta\)-可积的复值函数。假设下列条件之一成立：

\begin{enumerate}
\def\labelenumi{(\roman{enumi})}
\item
  \(f\) 与 \(\chi\) 连续；
\item
  G 在 X 中真作用，且 \(f\) 在一个可数紧集并集之外为零；
\item
  \(\mu\) 载于一个可数紧集并集。
\end{enumerate}

若 \(\mu\) 与 \(f\) 可作卷积，则函数 \(s \mapsto f(s^{-1}x)\chi(s^{-1}, x)\) 是本质 \(\mu\)-可积的，但一个局部 \(\beta\)-可忽略的 \(x\) 值集合除外，且局部 \(\beta\)-几乎处处有

\[
(\mu * ^ {\beta} f) (x) = \int_ {G} f (s ^ {- 1} x) \chi (s ^ {- 1}, x) d \mu (s) = \int_ {G} (\gamma_ {\chi} (s) f) (x) d \mu (s).\tag{3}
\]

设 \(h \in \mathcal{K}(\mathrm{X})\)。由于 \(\mu\) 与 \(f\) 可作卷积，函数 \((s, x) \mapsto h(sx)f(x)\) 是本质 \((\mu \otimes \beta)\)-可积的。由引理 1，函数 \((s, x) \mapsto h(x)f(s^{-1}x)\chi(s^{-1}, x)\) 是本质 \((\mu \otimes \beta)\)-可积的。在陈述的假设 (i) 或 (ii) 下，进而可推出该函数是 \((\mu \otimes \beta)\)-可积的；因为在第一种情形它连续，可应用第五章 §1, No.~1 命题 3，而在第二种情形它在一个可数紧集并集之外为零，可应用同上引 No.~2 命题 7, 2)。由 Lebesgue--Fubini 定理，

\[
\begin{array}{c} \iint h (s x) d \mu (s) d (f \cdot \beta) (x) = \iint h (x) f (s ^ {- 1} x) \chi (s ^ {- 1}, x) d \mu (s) d \beta (x) \\ = \int h (x) d \beta (x) \int f (s ^ {- 1} x) \chi (s ^ {- 1}, x) d \mu (s), \end{array}
\]

其中函数 \(x \mapsto g(x) = \int f(s^{-1}x)\chi (s^{-1},x)d\mu (s)\) 还是局部 \(\beta\)-可积的。于是可见

\[
\langle h, \mu * (f \cdot \beta) \rangle = \langle h, g \cdot \beta \rangle ,
\]

由此 \(g = \mu *^{\beta} f\)。

现在假设 \(\mu\) 载于由一列紧集组成的并集 \(S\)。函数

\[
(s, x) \mapsto h (x) f (s ^ {- 1} x) \chi (s ^ {- 1}, x) \varphi_ {\mathrm{S}} (s)
\]

是本质 \((\mu\otimes\beta)\)-可积的，且在一个可数紧集并集之外为零，从而是 \((\mu\otimes\beta)\)-可积的。由于 \(\mu=\varphi_{S}\cdot\mu\)，论证可如前一样完成。

备注. --- 当 \(\mu\) 有界时，命题 2 的假设 (iii) 特别地成立。因为，对每个 \(n > 0\)，存在 \(\mathbf{K}_n\)（\(\mathbf{G}\) 的一个紧子集），使得

\[
| \mu | (\mathrm{G} - \mathrm{K} _ {n}) \leqslant \frac {1}{n}
\]

（第四章 §4, No.~7），且 \(\mu\) 载于诸 \(K_{n}\) 的并集。更一般地，设 \(\rho\) 是一个 \(>0\) 的有限下半连续函数，定义在 \(G\) 上，满足 \(\rho(st) \leqslant \rho(s)\rho(t)\)；若 \(\mu \in \mathcal{M}^{\rho}\)，则假设 (iii) 成立；因为 \(\rho \cdot \mu\) 有界，且 \(\mu\) 与 \(\rho \cdot \mu\) 载于同样的子集，这是由于在 \(G\) 的每个紧子集上 \(\rho\) 都以一个 \(>0\) 的常数为下界。

\subsection{可作卷积的测度与函数的例子}

在命题 3 与命题 4 中，\(\mathcal{C}'(\mathrm{G})\) 与 \(\mathcal{M}(\mathrm{G})\) 分别赋予在 \(\mathcal{C}(\mathrm{G})\) 与 \(\mathcal{K}(\mathrm{G})\) 中的紧收敛拓扑。

命题 3. --- 假设 \(\chi\) 连续。设 \(\mu \in \mathcal{C}'(\mathbf{G})\)，\(f \in \mathcal{C}(\mathbf{X})\)。则：

\begin{enumerate}
\def\labelenumi{(\roman{enumi})}
\item
  \(\mu\) 与 \(f\) 关于 \(\beta\) 可作卷积。
\item
  No.~1 的公式 (3) 对每个 \(x \in \mathbf{X}\) 定义了一个连续的卷积 \(\mu *^{\beta}f\)，它正是元素 \(\pmb{\gamma}_{\chi}(\mu)f\)，即由连续表示 \(\pmb{\gamma}_{\chi}\)（\(\mathbf{G}\) 在 \(\mathcal{C}(\mathbf{X})\) 中）所定义者；此外，映射 \((\mu, f) \mapsto \mu *^{\beta}f\) 相对于 \(\mathcal{C}'(\mathbf{G})\) 的等度连续子集与 \(\mathcal{C}(\mathbf{X})\) 的紧子集是亚连续的。
\item
  若再有 \(f \in \mathcal{K}(\mathrm{X})\)，则 (ii) 中的乘积 \(\mu *^{\beta}f\) 属于 \(\mathcal{K}(\mathrm{X})\)，且映射 \((\mu, f) \mapsto \mu *^{\beta}f\) 相对于 \(\mathcal{C}'(\mathrm{G})\) 的等度连续子集与 \(\mathcal{K}(\mathrm{X})\) 的紧子集是亚连续的。
\end{enumerate}

我们知道 \(\mu\) 与 \(f\) 可作卷积（§3, No.~2 命题 8 (i)）。另一方面，采用 §2 的记号，我们有

\[
\boldsymbol {\gamma} _ {\chi} (\mu) f = \int \left(\boldsymbol {\gamma} _ {\chi} (s) f\right) d \mu (s) \in \mathcal {C} (\mathrm{X})
\]

因为 \(\mathcal{C}(\mathbf{X})\) 是完备的。特别地，对每个 \(x\in \mathbf{X}\)，

\[
\left(\boldsymbol {\gamma} _ {\chi} (\mu) f\right) (x) = \int \left(\boldsymbol {\gamma} _ {\chi} (s) f\right) (x) d \mu (s).
\]

把这一点与命题 2 (i) 以及 §2, No.~6 结合起来，便证明了 (ii)。最后，若 \(f \in \mathcal{K}(\mathrm{X})\)，则 \(\mu * (f \cdot \beta)\) 具有紧支集（§3, No.~2 命题 9），因此 \(\mu *^{\beta} f \in \mathcal{K}(\mathrm{X})\)。为此，考虑连续表示 \(U\)（\(\mathbf{G}\) 在完备化空间 \(\mathcal{K}(\mathrm{X})^{\wedge}\) 中的），它是由连续算子 \(\pmb{\gamma}_{\chi}(s)\)（在 \(\mathcal{K}(\mathbf{X})\) 中）连续延拓而得到的（§2, No.~1 备注 3）。设 S 为 \(\mu\) 的支集。诸函数 \(\pmb{\gamma}_{\chi}(s)f\)（\(s \in \mathbf{S}\)）的支集都包含在一个固定的紧集 K 中。集合 \(\mathcal{K}(\mathbf{X}, \mathbf{K})\) 是 \(\mathcal{K}(\mathbf{X})\) 的完备线性子空间。因此 \(U(\mu)f \in \mathcal{K}(\mathbf{X})\)。如前可见 \(U(\mu)f = \mu^{*\beta}f\)，而 (iii) 仍由 §2, No.~6 得出。

命题 4. --- 假设 G 在 X 中真作用且 \(\chi\) 连续。设 \(\mu \in \mathcal{M}(\mathrm{G})\)，\(f \in \mathcal{K}(\mathrm{X})\)。

\begin{enumerate}
\def\labelenumi{(\roman{enumi})}
\item
  \(\mu\) 与 \(f\) 关于 \(\beta\) 可作卷积。
\item
  No.~1 的公式 (3) 对每个 \(x \in \mathbf{X}\) 定义了一个连续的卷积 \(\mu *^{\beta}f\)。
\item
  映射 \((\mu, f) \mapsto \mu *^{\beta}f\)（从 \(\mathcal{M}(\mathrm{G}) \times \mathcal{K}(\mathrm{X})\) 到 \(\mathcal{C}(\mathrm{X})\) 中）相对于 \(\mathcal{M}(\mathrm{G})\) 的有界子集与 \(\mathcal{K}(\mathrm{X})\) 中包含在某个子空间 \(\mathcal{K}(\mathrm{X}, \mathrm{L})\) 内的紧子集（其中 \(\mathrm{L}\) 是 \(\mathrm{X}\) 的变动紧子集）是亚连续的。
\end{enumerate}

我们知道 \(\mu\) 与 \(f\) 可作卷积（§3, No.~2 命题 8 (ii)），且显然 (3) 中出现的积分对每个 \(x \in \mathbf{X}\) 都存在。设 \(\mathbf{K}\) 与 \(\mathbf{L}\) 是 \(\mathbf{X}\) 的两个紧子集。存在 \(\mathbf{H}\)（\(\mathbf{G}\) 的一个紧子集），使得关系 \(x \in \mathbf{K}\) 与 \(s^{-1}x \in \mathbf{L}\) 蕴含 \(s \in \mathbf{H}\)；取 \(\varphi \in \mathcal{H}_{+}(\mathbf{G})\) 满足 \(\varphi(s) = 1\)（对 \(s \in \mathbf{H}\)）。于是，对 \(f \in \mathcal{H}(\mathbf{X}, \mathbf{L})\) 与 \(x \in \mathbf{K}\)，

\[
\begin{array}{c} \int f (s ^ {- 1} x) \chi (s ^ {- 1}, x) d \mu (s) = \int f (s ^ {- 1} x) \chi (s ^ {- 1}, x) \varphi (s) d \mu (s) \\ = ((\varphi \cdot \mu) * ^ {\beta} f) (x). \end{array}
\]

因而 \(\int f(s^{-1}x)\chi(s^{-1},x)d\mu(s)\) 是 \(x\) 的连续函数，并定义了一个卷积 \(\mu * ^ {\beta} f\in \mathcal{C}(\mathrm{X})\)。此外，映射 \(\mu \mapsto \varphi \cdot \mu\)（从 \(\mathcal{M}(\mathrm{G})\) 到 \(\mathcal{C}'(\mathrm{G})\) 中）对于紧收敛拓扑是连续的。于是命题 3 (iii) 蕴含：映射 \((\mu ,f)\mapsto \mu * ^ {\beta} f\)（从 \(\mathcal{M}(\mathrm{G})\times \mathcal{K}(\mathrm{X},\mathrm{L})\) 到 \(\mathcal{C}(\mathrm{X})\) 中）对 X 的每个紧子集 L 相对于 \(\mathcal{K}(\mathrm{X},\mathrm{L})\) 的紧子集是亚连续的。特别地，映射 \((\mu ,f)\mapsto \mu * ^ {\beta} f\)（从 \(\mathcal{M}(\mathrm{G})\times \mathcal{K}(\mathrm{X})\) 到 \(\mathcal{C}(\mathrm{X})\) 中）是分别连续的。由于 \(\mathcal{K}(\mathrm{X})\) 是桶式空间，该映射相对于 \(\mathcal{M}(\mathrm{G})\) 的有界子集是亚连续的（TVS, III, §5, No.~3 命题 6）。

备注 1. --- 在命题 4 的假设下，映射 \(\mu \mapsto \mu *^{\beta}f\)（从 \(\mathcal{M}_{+}(\mathrm{G})\) 到 \(\mathcal{C}(\mathrm{X})\) 中）当 \(\mathcal{M}_{+}(\mathrm{G})\) 赋予淡拓扑时对每个 \(f\in \mathcal{K}(\mathrm{X})\) 是连续的。因为，设 K 为 X 的一个紧子集，S 为 \(f\) 的（紧）支集；由于 G 在 X 中真作用，由 \(s\in \mathrm{G}\) 中使得存在 \(x\in \mathrm{K}\) 满足 \(s^{-1}x\in \mathrm{S}\) 的那些 s 组成的集合是 G 的一个紧子集 L（GT, III, §4, No.~5 定理 1）。设 \(\varepsilon\) 为一个 \(>0\) 的数，\(\varphi\) 为 \(\mathcal{K}_{+}(\mathrm{G})\) 中在紧集 L 上等于 1 的函数，\(\mu_0\) 为 \(\mathcal{M}_{+}(\mathrm{G})\) 的一个元素；集合 \(\mathrm{W}_0\) 由测度 \(\mu \in \mathcal{M}_{+}(\mathrm{G})\) 组成，它们满足

\[
\left| \int \varphi (s) d \mu (s) - \int \varphi (s) d \mu_ {0} (s) \right| \leqslant \varepsilon
\]

是 \(\mu_0\) 在 \(\mathcal{M}_{+}(\mathrm{G})\) 中的一个邻域。另一方面，函数 \((s,x)\mapsto f(s^{-1}x)\chi (s^{-1},x)\) 在 \(\mathbf{L}\times \mathbf{K}\) 上一致连续，因此存在有限多个点 \(x_{i}\in \mathbf{K}\)（\(1\leqslant i\leqslant n\)），使得对每个 \(x\in \mathbf{K}\)，存在一个 \(i\)，使得

\[
\left| f (s ^ {- 1} x) \chi (s ^ {- 1}, x) - f (s ^ {- 1} x _ {i}) \chi (s ^ {- 1}, x _ {i}) \right| \leqslant \varepsilon
\]

对一切 \(s \in \mathbf{L}\) 成立。由于 \(\mu(\mathbf{L}) \leqslant \int \varphi(s) d\mu_0(s) + \varepsilon\) 对一切 \(\mu \in \mathbf{W}_0\) 成立，故还有

\[
\begin{array}{r l} \Big | \int f (s ^ {- 1} x) \chi (s ^ {- 1}, x) d \mu (s) - \int f (s ^ {- 1} x _ {i}) \chi (s ^ {- 1}, x _ {i}) d \mu (s) \Big | & \\ & \leqslant \varepsilon \Big (\int \varphi (s) d \mu_ {0} (s) + \varepsilon \Big) \end{array}
\]

对每个满足前一不等式的 x 以及每个 \(\mu \in W_{0}\) 成立。现在设 W 为 \(\mu_{0}\) 在 \(\mathcal{M}_{+}(G)\) 中的由测度 \(\mu \in W_{0}\) 组成的邻域，这些测度满足

\[
\left| \int f (s ^ {- 1} x _ {i}) \chi (s ^ {- 1}, x _ {i}) d \mu (s) - \int f (s ^ {- 1} x _ {i}) \chi (s ^ {- 1}, x _ {i}) d \mu_ {0} (s) \right| \leqslant \varepsilon
\]

对 \(1 \leqslant i \leqslant n\) 成立。显然，对每个测度 \(\mu \in \mathbf{W}\) 与每个 \(x \in \mathbf{K}\)，

\[
\begin{array}{r l} & {\Big | \int f (s ^ {- 1} x) \chi (s ^ {- 1}, x) d \mu (s) - \int f (s ^ {- 1} x) \chi (s ^ {- 1}, x) d \mu_ {0} (s) \Big |} \\ & {\qquad \leqslant \varepsilon \Big (2 \int \varphi (s) d \mu_ {0} (s) + 2 \varepsilon + 1 \Big),} \end{array}
\]

而由于 \(\varepsilon\) 是任意的，这就证明了我们的断言。

命题 5. --- 假设 \(\chi\) 是连续乘子且每个函数 \(\chi(s,\cdot)\) 有界。

\begin{enumerate}
\def\labelenumi{(\roman{enumi})}
\item
  函数 \(s \mapsto \rho(s) = \sup_{x \in X} \chi(s^{-1}, x)\)（定义在 \(G\) 上）是 \(> 0\) 的下半连续函数，且满足 \(\rho(st) \leqslant \rho(s)\rho(t)\)（对一切 \(s, t\)，\(G\) 中的）。
\item
  设 \(\mu \in \mathcal{M}^{\rho}(\mathrm{G})\) 且 \(f \in \mathrm{L}^{\infty}(\mathrm{X},\beta)\)。则 \(\mu\) 与 \(f\) 可作卷积，且 \(\mu *^{\beta} f\) 在局部几乎处处由 No.~1 的公式 (3) 给出。我们有 \(\mu *^{\beta} f \in \mathrm{L}^{\infty}(\mathrm{X},\beta)\)，且 \(\| \mu *^{\beta} f \|_{\infty} \leqslant \| \mu \|_{\rho} \| f \|_{\infty}\)。
\item
  若再有 \(f \in \mathcal{C}^{\infty}(\mathrm{X})\)（相应地，\(\overline{\mathcal{K}(\mathrm{X})}\)），则 No.~1 的公式 (3) 对每个 \(x\) 定义了一个卷积 \(\mu *^{\beta}f\)，它属于 \(\mathcal{C}^{\infty}(\mathrm{X})\)（相应地，\(\overline{\mathcal{K}(\mathrm{X})}\)）。
\item
  若 \(f \in \overline{\mathcal{K}(\mathrm{X})}\)，则由 (3) 定义的卷积 \(\mu *^{\beta}f\) 正是元素 \(\gamma_{\chi}(\mu)f\)，它由连续表示 \(\gamma_{\chi}\)（\(G\) 在 \(\overline{\mathcal{K}(\mathrm{X})}\) 中的）所定义。
\end{enumerate}

恒等式 \(\chi(st,x) = \chi(s,tx)\chi(t,x)\) 立即蕴含 \(\rho(st) \leqslant \rho(s)\rho(t)\)。另一方面，\(\rho\) 作为连续函数的上包络，是下半连续的。

设 \(\mu \in \mathcal{M}^{\rho}(\mathrm{G})\)。由 No.~1 的命题 1，\(\mu\) 与 1 可作卷积；命题 2 (i) 表明 \((|\mu|^{*\beta}1)(x) \leqslant \int_{\mathrm{G}} \rho(s)d|\mu|(s)\) 在局部 \(\beta\)-几乎处处成立。因此，若 \(f\) 是 \(\beta\)-可测的且 \(|f| \leqslant 1\)，则 \(\mu\) 与 \(f\) 可作卷积，且 \(\mathrm{N}_{\infty}(\mu^{*\beta}f) \leqslant \int \rho(s)d|\mu|(s)\)。此外，\(\mu^{*\beta}f\) 在局部几乎处处由 No.~1 的公式 (3) 给出，因为 No.~1 命题 2 的条件 (iii) 得以满足。这就蕴含 (ii)。

设 \(f\) 连续且绝对值以 1 为界。显然 (3) 中出现的积分对一切 \(x \in X\) 存在。我们来证明它们连续地依赖于 \(x\)。可以假设 \(\mu \geqslant 0\)。设 \(x_0 \in X\) 与 \(\varepsilon > 0\)。设 \(K\) 为 \(G\) 的一个紧子集，使得 \(\int_{G - K} \rho(s) d\mu(s) \leqslant \varepsilon\)。存在一个邻域 \(V\)（\(x_0\) 在 \(X\) 中的），使得 \(x \in V\) 蕴含

\[
\left| f \left(s ^ {- 1} x\right) \chi \left(s ^ {- 1}, x\right) - f \left(s ^ {- 1} x _ {0}\right) \chi \left(s ^ {- 1}, x _ {0}\right) \right| \leqslant \varepsilon / \mu (\mathrm{K})
\]

对 \(s \in \mathbf{K}\) 成立。于是，对 \(x \in \mathbf{V}\)，

\[
\begin{array}{r l} & {\Big | \int f (s ^ {- 1} x) \chi (s ^ {- 1}, x) d \mu (s) - \int f (s ^ {- 1} x _ {0}) \chi (s ^ {- 1}, x _ {0}) d \mu (s) \Big |} \\ & {\qquad \leqslant 2 \int_ {\mathrm{G-K}} \rho (s) d \mu (s) + \int_ {\mathrm{K}} \frac {\varepsilon}{\mu (\mathrm{K})} d \mu (s) \leqslant 3 \varepsilon ,} \end{array}
\]

由此得我们的断言。假设此外 \(f \in \overline{\mathcal{H}(\mathrm{X})}\)。设 \(\mathbf{H}\) 为 \(\mathbf{X}\) 的一个紧子集，使得 \(|f(y)| \leqslant \varepsilon\) 对 \(y \notin \mathbf{H}\) 成立。设 \(x \notin \mathbf{KH}\)。则 \(s^{-1}x \notin \mathbf{H}\) 对 \(s \in \mathbf{K}\) 成立，因此

\[
\begin{array}{r l} \Big | \int_ {\mathrm{G}} f (s ^ {- 1} x) \chi (s ^ {- 1}, x) d \mu (s) \Big | & \leqslant \int_ {\mathrm{G} - \mathrm{K}} \rho (s) d \mu (s) + \int_ {\mathrm{K}} \varepsilon \rho (s) d \mu (s) \\ & \leqslant \varepsilon \Big (1 + \int_ {\mathrm{G}} \rho (s) d \mu (s) \Big), \end{array}
\]

这就完成了 (iii) 的证明。

最后，若 \(f \in \overline{\mathcal{K}(\mathbf{X})}\)，则由于 \(\varepsilon_x \in \mathcal{M}^1(\mathbf{X})\) 对一切 \(x \in \mathbf{X}\) 成立，我们有

\[
\left(\boldsymbol {\gamma} _ {\chi} (\mu) f\right) (x) = \int \left(\boldsymbol {\gamma} _ {\chi} (s) f\right) (x) d \mu (s),
\]

从而 \(\gamma_{\chi}(\mu)f\) 就是由 (3) 定义的卷积 \(\mu *^{\beta}f\)。

命题 6. --- 假设 \(\chi\) 是连续乘子且每个函数 \(\chi(s, \cdot)\) 有界。设 \(\rho(s) = \sup_{x \in X} \chi(s^{-1}, x)\)。设 \(p\) 与 \(q\) 为一对共轭指数（\(1 \leqslant p < +\infty\)）。设 \(\mu \in \mathcal{M}^{\rho^{1/q}}(G)\) 且 \(f \in L^{p}(X, \beta)\)。则：

\begin{enumerate}
\def\labelenumi{(\roman{enumi})}
\item
  \(\mu\) 与 \(f\) 可作卷积；
\item
  卷积 \(\mu * ^ {\beta} f\) 在局部 \(\beta\)-几乎处处由公式 (3) 给出，且在局部 \(\beta\)-几乎处处等于一个函数 \(g\in \mathbf{L}^p (\mathbf{X},\beta)\)，使得 \(\| g\| _p\leqslant \| \mu \|_{\rho^{1 / q}}\| f\| _p;\) (iii) \(g\) 等于元素 \(\gamma_{\chi}(\mu)f\)，它由连续表示 \(\gamma_{\chi}\)（\(\mathbf{G}\) 在 \(\mathbf{L}^p (\mathbf{X},\beta)\) 中的）所定义。
\end{enumerate}

我们有

\[
\int^ {*} \| \gamma_ {\chi} (s) f \| _ {p} d | \mu | (s) \leqslant \left(\int^ {*} \rho (s) ^ {1 / q} d | \mu | (s)\right) \| f \| _ {p} <   + \infty
\]

（依据 §2, No.~5 的公式 (5)）。另一方面，映射 \(s \mapsto \pmb{\gamma}_{\chi}(s)f\)（从 G 到 \(\mathrm{L}^p (\mathbf{X},\beta)\) 中）是连续的（§2, No.~5 命题 9）。因此该映射是 \(\mu\)-可积的。设

\[
g = \int_ {\mathrm{G}} \left(\boldsymbol {\gamma} _ {\chi} (s) f\right) d \mu (s) \in \mathrm{L} ^ {p} (\mathrm{X}, \beta).
\]

我们有 \(\| g\| _p\leqslant \left(\int \rho^{1 / q}(s)d|\mu |(s)\right)\| f\| _p\)。把前面的备注应用于 \(|f|\)，可以看出映射 \(s\mapsto \varepsilon_s*\mid f|\)（从 G 到 \(\mathbf{L}^p (\mathbf{X},\beta)\) 中）是 \(\mu\)-可积的，因此对每个 \(h\in \mathcal{K}(\mathbf{X})\)，映射 \(s\mapsto \langle h,\varepsilon_s*(|f|\cdot \beta)\rangle\) 是 \(\mu\)-可积的。于是 §1, No.~5 的命题 7 证明 \(\mu\) 与 \(f\cdot \beta\) 可作卷积。此外，

\[
\begin{array}{r l} \int_ {\mathrm{X}} g (x) h (x) d \beta (x) & = \int_ {\mathrm{G}} d \mu (s) \int_ {\mathrm{X}} (\boldsymbol {\gamma} _ {\chi} (s) f) (x) h (x) d \beta (x) \\ & = \int_ {\mathrm{G}} \langle h, \varepsilon_ {s} * (f \cdot \beta) \rangle d \mu (s), \end{array}
\]

而由 §1, No.~5 的命题 7，最后这个积分等于 \(\langle h, \mu * (f \cdot \beta) \rangle\)。于是可见 g 是 \(\mu\) 与 f 的一个卷积。由命题 2 及其后的备注，这个卷积在局部 \(\beta\)-几乎处处由 (3) 给出。

按照记号的滥用，常常把陈述中的函数 g 之一记作 \(\mu * ^ {\beta} f\)，这使我们可以写出

\[
\left\| \mu * ^ {\beta} f \right\| _ {p} \leqslant \left\| \mu \right\| _ {\rho^ {1 / q}} \| f \| _ {p}.
\]

若 X 在无穷处可数，这种记法此外还是完全合理的。

推论. --- 在命题 6 的假设下，映射 \((\mu, f) \mapsto \mu *^{\beta} f\) 在 \(\mathrm{L}^{p}(\mathrm{X}, \beta)\) 上定义了 \(\mathcal{M}^{\rho^{1/q}}(\mathrm{G})\) 上的一个左模结构（\(1 \leqslant p \leqslant +\infty\)）。

这由命题 5、命题 6 以及卷积的结合性得出。

备注 2. --- 设 X 为局部紧空间，局部紧群 G 通过 \((x,s)\mapsto xs\) 连续地右作用于其上。设 \(\beta\) 为 X 上的正测度。设 \(\chi\) 为 \textgreater0 且定义在 \(G\times X\) 上的函数，它对 \(G\times X\) 上的每个测度可测，且 \(\boldsymbol{\delta}(s)\boldsymbol{\beta}=\chi(s,\cdot)\cdot\boldsymbol{\beta}\) 对每个 \(s\in G\) 成立。设 f 为 X 上局部 \(\beta\)-可积的函数，并设 \(\mu\) 为 G 上的测度。若 \(f\cdot\beta\) 与 \(\mu\) 可作卷积（对于映射 \((x,s)\mapsto xs\)，从 \(X\times G\) 到 X），则 \((f\cdot\beta)*\mu\) 以 \(\beta\) 为基。这时称 f 与 \(\mu\) 关于 \(\beta\) 可作卷积；\((f\cdot\beta)*\mu\) 关于 \(\beta\) 的每个密度都称为 f 与 \(\mu\) 关于 \(\beta\) 的一个卷积，记作 \(f * ^ {\beta} \mu\) 或简作 \(f*\mu\)。

设 \(G^{0}\) 为 G 的反群。通过 \((s,x)\mapsto xs\)，\(G^{0}\) 连续地左作用于 X。说 f 与 \(\mu\) 按前述意义可作卷积，等价于说 \(\mu\) 与 f 对于左作用于 X 的 \(G^{0}\) 可作卷积；而且卷积 \(f * ^ {\beta} \mu\) 不是别的，正是 \(\mu * ^ {\beta} f\) 对于左作用于 X 的 \(G^{0}\) 的卷积。另一方面，对 \(s\in G^{0}\) 有 \(\pmb{\gamma}(s)\beta = \chi(s^{-1},\cdot)\cdot \beta\)。于是 No.~1 与 No.~2 的结果可以立即翻译为关于乘积 \(f * ^ {\beta} \mu\) 的结果。特别地：

\begin{enumerate}
\def\labelenumi{\arabic{enumi})}
\tightlist
\item
  若 \(s \in \mathbf{G}\) 且 \(f\) 局部 \(\beta\)-可积，则 \(f\) 与 \(\varepsilon_s\) 可作卷积，且
\end{enumerate}

\[
(f * \varepsilon_ {s}) (x) = \chi (s ^ {- 1}, x) f (x s ^ {- 1}).\tag{4}
\]

\begin{enumerate}
\def\labelenumi{\arabic{enumi})}
\setcounter{enumi}{1}
\tightlist
\item
  若 \(f\) 与 \(\mu\) 可作卷积，且命题 2 的条件 (i)、(ii)、(iii) 之一成立，则 \(f *^{\beta} \mu\) 在局部 \(\beta\)-几乎处处由下式给出
\end{enumerate}

\[
(f * ^ {\beta} \mu) (x) = \int_ {G} f (x s ^ {- 1}) \chi (s ^ {- 1}, x) d \mu (s).\tag{5}
\]

我们把翻译其余陈述的任务留给读者。注意，若 \(\chi\) 连续且 \(\beta\) 的支集为 X，则

\[
\chi (t s, x) = \chi (s, x t) \chi (t, x) \quad (x \in \mathrm{X}; s, t \text {in} \mathrm{G}).\tag{6}
\]

\subsection{卷积与转置}

我们回到 No.~1 开头处的假设与记号，但此外假设 \(\beta\) 是以 \(\chi\) 为乘子的相对不变测度；于是 \(\chi\) 是 G 上的连续函数。

命题 7. --- 设 \(f\) 是局部 \(\beta\)-可积的函数，定义在 \(X\) 上，\(\nu\) 是 \(X\) 上的测度，\(\mu\) 是 \(G\) 上的测度。假设：

\begin{enumerate}
\def\labelenumi{(\roman{enumi})}
\item
  \(\mu\) 与 \(f\) 可作卷积，且 No.~1 的公式 (3) 在局部 \(\beta\)-几乎处处定义了一个卷积 \(\mu *^{\beta} f\)。
\item
  \(\chi \cdot \check{\mu}\) 与 \(\nu\) 可作卷积。
\item
  函数 \(g(s,x) = f(s^{-1}x)\chi (s^{-1})\) 是 \((\mu \otimes \nu)\)-可积的。
\end{enumerate}

则 \(f\) 对 \((\chi \cdot \breve{\mu}) * \nu\) 是本质可积的，由 (3) 定义的函数 \(\mu *^{\beta} f\) 是 \(\nu\)-可积的，且

\[
\nu (\mu * ^ {\beta} f) = \big ((\chi \cdot \check {\mu}) * \nu \big) (f).\tag{7}
\]

由于 \(g(s, x)\) 对 \(\mu \otimes \nu\) 可积，函数 \(f(sx)\) 对 \((\chi \cdot \breve{\mu}) \otimes \nu\) 是本质可积的，且 \(f\) 对 \((\chi \cdot \breve{\mu}) * \nu\) 是本质可积的。由 Lebesgue--Fubini 定理，\(\mu *^{\beta} f = \int g(s, x) d\mu(s)\) 是 \(\nu\)-可积的，且

\[
\begin{array}{l} \nu (\mu * ^ {\beta} f) = \iint f (s ^ {- 1} x) \chi (s ^ {- 1}) d \mu (s) d \nu (x) \\ = \iint f (s x) \chi (s) d \stackrel {\vee} {\mu} (s) d \nu (x) = \big ((\chi \cdot \stackrel {\vee} {\mu}) * \nu \big) (f). \end{array}
\]

例. --- 1) 由命题 3 以及 §1, No.~4 命题 5 的推论，可取 \(f \in \mathcal{C}(\mathrm{X})\)，\(\nu \in \mathcal{C}'(\mathrm{X})\) 与 \(\mu \in \mathcal{C}'(\mathrm{G})\)。这时公式 (7) 表示自同态 \(\nu \mapsto (\chi \cdot \breve{\mu}) * \nu\)（\(\mathcal{C}'(\mathrm{X})\) 的）是自同态 \(f \mapsto \mu * f\)（\(\mathcal{C}(\mathrm{X})\) 的）的转置。

\begin{enumerate}
\def\labelenumi{\arabic{enumi})}
\setcounter{enumi}{1}
\item
  可取 \(f \in \mathcal{K}(\mathrm{X})\)，\(\nu \in \mathcal{M}(\mathrm{X})\) 与 \(\mu \in \mathcal{C}'(\mathrm{G})\)（依据命题 3、§3, No.~2 的命题 8，以及连续函数 \(g(s,x)\) 的支集与 \(\mu \otimes \nu\) 的支集相交于一个紧集这一备注）。这时公式 (7) 表示自同态 \(\nu \mapsto (\chi \cdot \breve{\mu}) * \nu\)（\(\mathcal{M}(\mathrm{X})\) 的）是自同态 \(f \mapsto \mu * f\)（\(\mathcal{K}(\mathrm{X})\) 的）的转置。
\item
  若 G 真作用于 X，则可取 \(f \in \mathcal{K}(X)\)，\(\nu \in \mathcal{C}'(X)\) 与 \(\mu \in \mathcal{M}(G)\)（依据命题 4、§3, No.~2 的命题 8 以及例 2 中的同样备注）。
\end{enumerate}

命题 8. --- 设 \(f\) 与 \(g\) 是两个局部 \(\beta\)-可积的函数，定义在 \(X\) 上，并设 \(\mu \in \mathcal{M}(G)\)。假设：

\begin{enumerate}
\def\labelenumi{(\roman{enumi})}
\item
  \(\mu\) 与 \(f\) 可作卷积，且 No.~1 的公式 (3) 在局部 \(\beta\)-几乎处处定义了一个卷积 \(\mu *^{\beta} f\)。
\item
  \(\chi \cdot \breve{\mu}\) 与 \(g\) 可作卷积，且 No.~1 的公式 (3)（把 \(\mu\) 换成 \(\chi \cdot \breve{\mu}\)、把 \(f\) 换成 \(g\)）在局部 \(\beta\)-几乎处处定义了一个卷积 \((\chi \cdot \breve{\mu}) *^{\beta} g\)。
\item
  存在一个函数 \(\psi\)（定义在 \(\mathbf{G}\) 上），它局部 \(\mu\)-几乎处处等于 1，使得函数
\end{enumerate}

\[
h (s, x) = g (x) f \left(s ^ {- 1} x\right) \chi \left(s ^ {- 1}\right) \psi (s)
\]

是 \((\mu \otimes \beta)\)-可积的。

则函数 \(g(x)\big((\mu * ^ {\beta} f)(x)\big)\) 与 \(f(x)\big(\big((\chi \cdot \breve{\mu}) * ^ {\beta} g\big)(x)\big)\) 是本质 \(\beta\)-可积的，且

\[
\int f (x) \big (\big ((\chi \cdot \breve {\mu}) * ^ {\beta} g \big) (x) \big) d \beta (x) = \int g (x) \big ((\mu * ^ {\beta} f) (x) \big) d \beta (x).\tag{8}
\]

因为，由 (iii) 与 Lebesgue--Fubini 定理，函数

\[
x \mapsto g (x) \int f (s ^ {- 1} x) \chi (s ^ {- 1}) \psi (s) d \mu (s)
\]

是 \(\beta\)-可积的，且

\[
\begin{array}{l} \mathrm{I} = \iint f (s ^ {- 1} x) g (x) \chi (s ^ {- 1}) \psi (s) d \mu (s) d \beta (x) \\ = \int g (x) d \beta (x) \int f (s ^ {- 1} x) \chi (s ^ {- 1}) \psi (s) d \mu (s). \end{array}
\]

但 \(\psi \cdot \mu = \mu\)，因此

\[
\int f (s ^ {- 1} x) \chi (s ^ {- 1}) \psi (s) d \mu (s) = (\mu * ^ {\beta} f) (x)
\]

在局部 \(\beta\)-几乎处处成立。这表明函数

\[
x \mapsto g (x) \left(\left(\mu * ^ {\beta} f\right) (x)\right)
\]

是本质 \(\beta\)-可积的，且

\[
\mathrm{I} = \int g (x) \left(\left(\mu * ^ {\beta} f\right) (x)\right) d \beta (x).
\]

另一方面，引理 1 表明函数

\[
(s, x) \mapsto g (s x) f (x) \chi (s ^ {- 1}) \psi (s)
\]

对 \((\chi \cdot \mu) \otimes \beta\) 可积。因此函数 \((s, x) \mapsto g(s^{-1}x)f(x)\psi(s^{-1})\) 对 \(\breve{\mu} \otimes \beta\) 可积，且

\[
\begin{array}{l} \mathrm{I} = \iint g (s ^ {- 1} x) f (x) \psi (s ^ {- 1}) d \mu^ {\vee} (s) d \beta (x) \\ = \int f (x) d \beta (x) \int g (s ^ {- 1} x) \psi (s ^ {- 1}) d \mu^ {\vee} (s). \end{array}
\]

但 \(\check{\psi} \cdot \check{\mu} = \check{\mu}\)，于是 \(\int g(s^{-1}x)\psi(s^{-1})d\check{\mu}(s) = ((\chi \cdot \check{\mu}) *^{\beta}g)(x)\) 在局部 \(\beta\)-几乎处处成立。这表明函数

\[
x \mapsto f (x) \left(\left((\chi \cdot \stackrel {\vee} {\mu}) * ^ {\beta} g\right) (x)\right)
\]

是本质 \(\beta\)-可积的，且

\[
\mathrm{I} = \int f (x) \left(\left(\left(\chi \cdot \stackrel {\vee} {\mu}\right) * ^ {\beta} g\right) (x)\right) d \beta (x).
\]

这就证明了命题。

例. --- 4) 可取 \(f \in \mathcal{C}(\mathrm{X})\)，\(g \in \mathcal{K}(\mathrm{X})\) 与 \(\mu \in \mathcal{C}'(\mathrm{G})\)（取 \(\psi = 1\)）。

\begin{enumerate}
\def\labelenumi{\arabic{enumi})}
\setcounter{enumi}{4}
\item
  若 \(\mathbf{G}\) 真作用于 \(\mathbf{X}\)，则可取 \(f \in \mathcal{K}(\mathbf{X})\)，\(g \in \mathcal{K}(\mathbf{X})\) 与 \(\mu \in \mathcal{M}(\mathbf{G})\)（取 \(\psi = 1\)）。
\item
  可取 \(f \in \mathrm{L}^p(\mathrm{X},\beta)\)，\(g \in \mathrm{L}^q(\mathrm{X},\beta)\) 与 \(\mu \in \mathcal{M}^\rho(\mathrm{G})\)，其中 \(1 \leqslant p < +\infty\)，\(\frac{1}{p} + \frac{1}{q} = 1\)，\(\rho = \chi^{-1/q}\)。条件 (i) 与 (ii) 由命题 5 与命题 6 满足。我们来证明 (iii)。我们已经看到，\(\mu\) 载于一个集合 \(S\)，它是一列紧集的并。取 \(\psi\) 为 \(S\) 的特征函数。函数 \(h\) 是 \((\mu \otimes \beta)\)-可测的：因为函数 \((s,x) \mapsto g(x)\chi(s^{-1})\psi(s)\) 如此，且由引理 1，函数 \((s,x) \mapsto f(s^{-1}x)\) 也如此。此外，由于 \(g\) 在一列 \(\beta\)-可积集合的并之外为零，\(h\) 在一列 \((\mu \otimes \beta)\)-可积集合的并之外为零。于是我们有（第五章 §8, No.~3 命题 7）：
\end{enumerate}

\[
\begin{array}{l} \mathrm{J} = \iint^ {*} | g (x) f (s ^ {- 1} x) | \chi (s ^ {- 1}) \psi (s) d | \mu | (s) d \beta (x) \\ = \int^ {*} | g (x) | d \beta (x) \int^ {*} | f (s ^ {- 1} x) | \chi (s ^ {- 1}) \psi (s) d | \mu | (s). \end{array}\tag{9}
\]

而由于 \(g\)（相应地，\(\psi\)）在一列可积集合的并之外为零，(9) 式右端的上积分等于本质上积分（第五章 §1, No.~2 命题 7）。现在（第五章 §5, No.~3 命题 3）

\[
\int^ {\bullet} | f (s ^ {- 1} x) | \chi (s ^ {- 1}) \psi (s) d | \mu | (s) = \int^ {\bullet} | f (s ^ {- 1} x) | \chi (s ^ {- 1}) d | \mu | (s)
\]

由于 \(\mu = \psi \cdot \mu\)。由命题 6，最后这个积分有限，且等于 \((|\mu| *^{\beta} |f|)(x)\)（在局部 \(\beta\)-几乎处处）。因此

\[
J = \int^ {\bullet} | g (x) | \left(| \mu | * ^ {\beta} | f |\right) (x) d \beta (x),
\]

而 J 有限，因为 \(g \in \mathbf{L}^q\) 且 \(|\mu| *^\beta |f| \in \mathbf{L}^p\)（命题 6）。因此 \(h\) 是 \((\mu \otimes \beta)\)-可积的。

这时公式 (8) 表示自同态 \(g \mapsto (\chi \cdot \breve{\mu}) * g\)（\(\mathrm{L}^q(\mathrm{X},\beta)\) 的）对 \(\mu \in \mathcal{M}^\rho(\mathrm{G})\) 是自同态 \(f \mapsto \mu * f\)（\(\mathrm{L}^p(\mathrm{X},\beta)\) 的）的转置。

\subsection{群上一个测度与一个函数的卷积}

设 G 为局部紧群。在 No.~4 与 No.~5 中，我们固定 G 上一个 \(\beta \neq 0\) 的相对不变正测度；设 \(\chi\) 与 \(\chi'\) 分别是它的左乘子与右乘子（回忆 \(\chi' = \chi \Delta_{\mathrm{G}}\)）。若 \(\mu\) 是 G 上的测度而 \(f\) 是 G 上局部 \(\beta\)-可积的函数，则 \(\mu\) 与 \(f\) 的可卷积性以及乘积 \(\mu * f\)（相应地，\(f\) 与 \(\mu\) 的可卷积性以及乘积 \(f * \mu\)）可通过把 G 视为经平移左作用于自身（相应地，右作用于自身）来定义。让我们把前面的一些结果在这一情形下明确写出：

\begin{enumerate}
\def\labelenumi{\arabic{enumi})}
\tightlist
\item
  设 \(\mu\) 为 \(\mathbf{G}\) 上的测度，\(f\) 为一个局部 \(\beta\)-可积的复值函数（\(\mathbf{G}\) 上的）。假设下列条件之一已被验证：
\end{enumerate}

\begin{enumerate}
\def\labelenumi{(\roman{enumi})}
\item
  \(f\) 连续；
\item
  \(f\) 在一列紧集的并的余集上为零；
\item
  \(\mu\) 载于一列紧集的并。
\end{enumerate}

若 \(\mu\) 与 \(f\) 可作卷积，则在局部 \(\beta\)-几乎处处，

\[
(\mu * f) (x) = \int_ {G} f (s ^ {- 1} x) \chi (s ^ {- 1}) d \mu (s).\tag{10}
\]

若 \(f\) 与 \(\mu\) 可作卷积，则在局部 \(\beta\)-几乎处处，

\[
(f * \mu) (x) = \int_ {\mathrm{G}} f (x s ^ {- 1}) \chi^ {\prime} (s ^ {- 1}) d \mu (s).\tag{11}
\]

\begin{enumerate}
\def\labelenumi{\arabic{enumi})}
\setcounter{enumi}{1}
\tightlist
\item
  设 \(p\) 与 \(q\) 为两个共轭指数（\(1 \leqslant p \leqslant +\infty\)）。若 \(\mu \in \mathcal{M}^{\chi^{-1/q}}(\mathrm{G})\) 且 \(f \in \mathrm{L}^p(\mathrm{G},\beta)\)，则 \(\mu\) 与 \(f\) 可作卷积，且 \(\mu * f\) 在局部 \(\beta\)-几乎处处等于 \(\mathrm{L}^p(\mathrm{G},\beta)\) 中的一个函数；我们有（采用已经指出的记号滥用）
\end{enumerate}

\[
\| \mu * f \| _ {p} \leqslant \| \mu \| _ {\chi^ {- 1 / q}} \| f \| _ {p}.
\]

若 \(\mu \in \mathcal{M}^{\chi'^{-1/q}}(\mathrm{G})\) 且 \(f \in \mathrm{L}^p(\mathrm{G},\beta)\)，则 \(f\) 与 \(\mu\) 可作卷积，且 \(f * \mu\) 在局部 \(\beta\)-几乎处处等于 \(\mathrm{L}^p(\mathrm{G},\beta)\) 中的一个函数；我们有 \(\| f * \mu \|_p \leqslant \| \mu \|_{\chi'^{-1/q}} \| f \|_p\)。

\begin{enumerate}
\def\labelenumi{\arabic{enumi})}
\setcounter{enumi}{2}
\item
  映射 \((\mu, f) \mapsto \mu * f\) 与 \((f, \mu) \mapsto f * \mu\) 在 \(\mathbf{L}^p(\mathbf{G}, \beta)\) 上定义了 \(\mathcal{M}^{\chi^{-1/q}}(\mathbf{G})\) 上的左模结构与 \(\mathcal{M}^{\chi'^{-1/q}}(\mathbf{G})\) 上的右模结构。\(\mathbf{L}^p(\mathbf{G}, \beta)\) 上的这两个外运算由卷积的结合性而可交换。
\item
  若 \(\mu * f\) 连续且在每一点由 (10) 给出，则
\end{enumerate}

\[
(\mu * f) (e) = \int f (s ^ {- 1}) \chi (s ^ {- 1}) d \mu (s).\tag{12}
\]

若 \(f * \mu\) 连续且在每一点由 (11) 给出，则

\[
(f * \mu) (e) = \int f (s ^ {- 1}) \chi^ {\prime} (s ^ {- 1}) d \mu (s).\tag{13}
\]

\subsection{群上函数的卷积}

我们保留 No.~4 的记号 G、\(\beta\)、\(\chi\)、\(\chi'\)。

回忆：若 f 是 G 上的复值函数，则局部 \(\beta\)-可积这一性质与 \(\beta\) 的选取无关。设 \(\mathcal{L}(\mathrm{G})\) 为具有这一性质的函数组成的集合。若 \(f \in \mathcal{L}(\mathrm{G})\)，\(g \in \mathcal{L}(\mathrm{G})\)，则关系

\[\text{«}f \cdot \beta\text{ 与 }g \cdot \beta\text{ 可作卷积 »}\]

与 \(\beta\) 的选取无关（§3, No.~1 命题 6）。这时我们就说 \(f\) 与 \(g\) 可作卷积。由 No.~1，\((f\cdot \beta)*(g\cdot \beta)\) 具有 \(h\cdot \beta\) 的形式，其中 \(h\in \mathcal{L}(\mathbf{G})\)，\(h\) 在局部 \(\beta\)-可忽略集的范围内被确定。我们将记 \(h = f * ^ {\beta} g\)，并说 \(h\) 是 \(f\) 与 \(g\) 相对于 \(\beta\) 的一个卷积。（当不致混淆时省略 \(\beta\)。）若把 \(\beta\) 换成 \(\psi \cdot \beta\)，其中 \(\psi\) 是 \(\mathbf{G}\) 在 \(\mathbf{R}_+^*\) 中的一个连续表示，则 \(h\) 不变（§3, No.~1 命题 6）；若把 \(\beta\) 换成 \(a\beta\)（\(a\in \mathbf{R}_+^*\)），则 \(h\) 换成 ah。G 上若干个函数的卷积以类似方式定义。

若 f 与 g 的卷积之一连续，则它被唯一确定，因为 \(\beta\) 的支集是 G。这时称它为 f 与 g 相对于 \(\beta\) 的卷积。

显然

\[
f * ^ {\beta} g = (f \cdot \beta) * ^ {\beta} g = f * ^ {\beta} (g \cdot \beta).\tag{14}
\]

命题 9. --- 设 \(f,g\) 属于 \(\mathcal{L}(\mathbf{G})\)。假设函数 \(s\mapsto g(s^{-1}x)f(s)\chi(s^{-1})\) 是本质 \(\beta\)-可积的，除去一个局部 \(\beta\)-可忽略的 \(x\) 值集合，并且函数

\[
x \mapsto \int | g (s ^ {- 1} x) f (s) | \chi (s ^ {- 1}) d \beta (s),
\]

它局部 \(\beta\)-几乎处处有定义且局部 \(\beta\)-可积。则 \(f\) 与 \(g\) 可作卷积。

这由 No.~1 的命题 1 得出。

命题 10. --- 设 \(f, g\) 属于 \(\mathcal{L}(G)\)。假设这两个函数之一连续，或在一列紧集的并的余集上为零。若 \(f\) 与 \(g\) 可作卷积，则函数 \(f * g\) 在局部 \(\beta\)-几乎处处由下式给出

\[
\begin{array}{c} (f * g) (x) = \int_ {\mathrm{G}} g (s ^ {- 1} x) f (s) \chi (s ^ {- 1}) d \beta (s) \\ = \int_ {\mathrm{G}} f (x s ^ {- 1}) g (s) \chi^ {\prime} (s ^ {- 1}) d \beta (s). \end{array}\tag{15}
\]

这由 No.~1 的命题 2 以及 No.~4 中的备注得出。

特别地，若 \(f * g\) 连续且在每一点由 (15) 给出，则

\[
(f * g) (e) = \int g (s ^ {- 1}) f (s) \chi (s ^ {- 1}) d \beta (s) = \int f (s ^ {- 1}) g (s) \chi^ {\prime} (s ^ {- 1}) d \beta (s).\tag{16}
\]

更特别地，若 \(\beta\) 是左 Haar 测度且是右 Haar 测度，且 \(f*g\) 与 \(g*f\) 连续并在每一点由 (15) 与 \(g*f\) 的类似公式给出，则

\[
(f * g) (e) = (g * f) (e) = \int f (s) g (s ^ {- 1}) d \beta (s).\tag{17}
\]

命题 11. --- 设 \(f,g\) 属于 \(\mathcal{L}(\mathrm{G})\)。假设函数 \(f,g\) 之一连续，且函数 \(f,g\) 之一具有紧支集。则 \(f\) 与 \(g\) 可作卷积。公式 (15) 对所有 \(x\in\mathrm{G}\) 定义了一个连续的乘积 \(f*g\)。若 \(f\in\mathcal{K}(\mathrm{G})\) 且 \(g\in\mathcal{K}(\mathrm{G})\)，则 \(f*g\in\mathcal{K}(\mathrm{G})\)。

这由 No.~2 的命题 3 与命题 4 得出。

命题 12. --- 设 \(p\) 与 \(q\) 为两个共轭指数 \((1 \leqslant p \leqslant +\infty)\)。若 \(f\chi^{-1/q} \in \mathrm{L}^1(\mathrm{G}, \beta)\) 且 \(g \in \mathrm{L}^p(\mathrm{G}, \beta)\)，则 \(f\) 与 \(g\) 可作卷积，\(f * g\) 在局部 \(\beta\)-几乎处处等于 \(\mathrm{L}^p(\mathrm{G}, \beta)\) 中的一个函数，且

\[
\left\| f * g \right\| _ {p} \leqslant \left\| f \chi^ {- 1 / q} \right\| _ {1} \left\| g \right\| _ {p}.
\]

若 \(f \in \mathrm{L}^p(\mathbf{G},\beta)\) 且 \(g\chi'^{-1/q} \in \mathrm{L}^1(\mathbf{G},\beta)\)，则 \(f\) 与 \(g\) 可作卷积，\(f * g\) 在局部 \(\beta\)-几乎处处等于 \(\mathrm{L}^p(\mathbf{G},\beta)\) 中的一个函数，且

\[
\| f * g \| _ {p} \leqslant \| f \| _ {p} \| g \chi^ {\prime - 1 / q} \| _ {1}.
\]

这由 No.~2 的命题 5 与命题 6 以及 No.~4 中的备注得出。

命题 13. --- 若 \(f\chi^{-1} \in \mathbf{L}^1(\mathbf{G}, \beta)\) 且 \(g \in \overline{\mathcal{H}(\mathbf{G})}\)，或者若 \(f \in \overline{\mathcal{H}(\mathbf{G})}\) 且 \(g\chi'^{-1} \in \mathbf{L}^1(\mathbf{G}, \beta)\)，则 \(f\) 与 \(g\) 可作卷积，且 (15) 对每一 \(x \in \mathbf{G}\) 定义了一个乘积 \(f * g\)，它属于 \(\overline{\mathcal{H}(\mathbf{G})}\)。

这由 No.~2 的命题 5 以及 No.~4 中的备注得出。

命题 14. --- 若 \(f\chi^{-1} \in \mathbf{L}^1(\mathbf{G}, \beta)\) 且 \(g \in \mathbf{L}^\infty(\mathbf{G}, \beta)\)，则公式 (15) 对每一 \(x \in \mathbf{G}\) 定义了一个有界的乘积 \(f * g\)，它对 \(\mathbf{G}\) 的右一致结构是一致连续的。

我们已经知道 \(f * g\) 属于 \(\mathrm{L}^{\infty}(\mathrm{G},\beta)\)（No.~2 命题 5）；此外，\((f*g)(x) = \int f(xs^{-1})g(s)d\nu(s)\)，其中置 \(\nu = \chi'^{-1}\cdot \beta\)；\(\nu\) 是一个右 Haar 测度。因此

\[
\begin{array}{c} | (f * g) (x) - (f * g) (x ^ {\prime}) | \leqslant \| g \| _ {\infty} \int | f (x s ^ {- 1}) - f (x ^ {\prime} s ^ {- 1}) | d \nu (s) \\ = \| g \| _ {\infty} \int | (f (s ^ {- 1}) - f (x ^ {\prime} x ^ {- 1} s ^ {- 1}) d \nu (s) \end{array}
\]

且后一积分可以任意小，只要 \(x'x^{-1}\) 属于 \(e\) 的一个适当邻域（§2, No.~5 命题 8）。

命题 15. --- 设 \(p\) 与 \(q\) 为两个共轭指数 \((1 < p < +\infty)\)。假设 \(\beta\) 是左不变的。设 \(f \in \mathrm{L}^p(\mathrm{G}, \beta)\)，\(g \in \mathrm{L}^q(\mathrm{G}, \breve{\beta})\)。则 \(f\) 与 \(g\) 可作卷积。公式 (15) 对每一 \(x \in \mathrm{G}\) 定义了一个乘积 \(f * g\)，它属于 \(\overline{\mathcal{H}(\mathrm{G})}\) 且满足 \(\| f * g \|_{\infty} \leqslant \| f \|_p \| \breve{g} \|_q\)。

因为，我们有 \(\check{g} \in \mathbf{L}^q(\mathbf{G}, \beta)\)，因此函数 \(s \mapsto g(s^{-1}x)f(s)\) 是 \(\beta\)-可积的，对每一 \(x \in \mathbf{G}\)。此外，

\[
\begin{array}{c} \int | g (s ^ {- 1} x) f (s) | d \beta (s) \leqslant \bigg (\int | f (s) | ^ {p} d \beta (s) \bigg) ^ {1 / p} \bigg (\int | g (s ^ {- 1} x) | ^ {q} d \beta (s) \bigg) ^ {1 / q} \\ = \| f \| _ {p} \bigg (\int | \check {g} (x ^ {- 1} s) | ^ {q} d \beta (s) \bigg) ^ {1 / q} = \| f \| _ {p} \| \check {g} \| _ {q}, \end{array}
\]

因此 \(f\) 与 \(g\) 可作卷积（命题 9）。同时可以看到 (15) 对每一 \(x\) 定义了一个乘积 \(f * g\)，使得

\[
\left| (f * g) (x) \right| \leqslant \| f \| _ {p} \| \check {g} \| _ {q}.
\]

对 \(f, g\)（\(\mathcal{H}(\mathrm{G})\) 中的），我们有 \(f * g \in \mathcal{H}(\mathrm{G})\)（命题 11）；因此，对 \(f \in \mathrm{L}^p(\mathrm{G}, \beta)\) 与 \(g \in \mathrm{L}^q(\mathrm{G}, \beta)\)，由 (15) 给出的乘积 \(f * g\) 是 \(\mathcal{H}(\mathrm{G})\) 中函数的一致极限，故属于 \(\overline{\mathcal{H}(\mathrm{G})}\)。

推论. --- 设 \(f \in \mathbf{L}^2(\mathbf{G}, \beta)\)，\(g \in \mathbf{L}^2(\mathbf{G}, \beta)\)。则 \(f\) 与 \(\breve{g}\) 可作卷积。卷积 \(f * \breve{g}\) 之一属于 \(\overline{\mathcal{H}(\mathbf{G})}\)，且它在 \(e\) 处的值为 \(\int_{\mathbf{G}} f(s) \overline{g(s)} d\beta(s)\)。

只需在命题 15 中取 \(p = q = 2\) 并应用 (16)。

我们不再假设 \(\beta\) 是左不变的。设 \(\rho\) 为 G 上的一个下半连续的有限函数 \textgreater0，使得对 G 中所有 s, t 有 \(\rho(st) \leqslant \rho(s)\rho(t)\)。我们用 \(\mathrm{L}^{\rho}(\mathrm{G}, \beta)\) 表示 G 上对 \(\rho \cdot \beta\) 可积的复值函数的等价类组成的集合。借助映射 \(f \mapsto f \cdot \beta\)，\(\mathrm{L}^{\rho}(\mathrm{G}, \beta)\) 可等同于 \(\mathcal{M}^{\rho}(\mathrm{G})\) 中以 \(\beta\) 为基的元素组成的集合（该集合与 \(\beta\) 的选取无关）。若置

\[
\| f \| _ {\rho} = \int_ {G} | f (s) | \rho (s) d \beta (s)
\]

对 \(f \in \mathrm{L}^{\rho}(\mathrm{G},\beta)\)，这一等同与范数相容，于是 \(\mathrm{L}^{\rho}(\mathrm{G},\beta)\) 表现为 \(\mathcal{M}^{\rho}(\mathrm{G})\) 的一个完备赋范子代数。由 §3, No.~2 命题 10，它甚至是 \(\mathcal{M}^{\rho}(\mathrm{G})\) 的一个双边理想。（取 \(\rho = 1\)，就重新得到 No.~4 的断言之一。）特别地，\(\mathrm{L}^{1}(\mathrm{G},\beta)\) 可等同于 \(\mathcal{M}^{1}(\mathrm{G})\) 的一个闭双边理想。

命题 16. --- 设 U 为 G 在 Banach 空间 E 中的一个连续表示。置 \(\rho(s) = \|U(s)\|\)（对所有 \(s \in G\)）。对每一 \(f \in L^{\rho}(G, \beta)\)，置 \(U(f) = U(f \cdot \beta)\)。则 \(f \mapsto U(f)\) 是代数 \(L^{\rho}(G, \beta)\) 在 E 中的一个线性表示，满足 \(\|U(f)\| \leqslant \|f\|_{\rho}\)。

这由 §2, No.~6 与 §3, No.~3 命题 11 得出。

\subsection{应用}

命题 17. --- 设 G 为局部紧群，A 为 G 的一个子集，它对一个 Haar 测度可测且非局部可忽略。则 \(A \cdot A^{-1}\) 是 e 的一个邻域。

设 \(\beta\) 为一个左 Haar 测度。存在紧子集 \(K\)（\(G\) 的），使得 \(B = A \cap K\) 可积，测度 \(>0\)（对 \(\beta\)）。我们对 \(f = g = \varphi_{B}\) 应用命题 15 的推论。函数 \(F = \varphi_{B} * \check{\varphi}_{B}\) 连续且 \(>0\)（在 \(e\) 处）。因此存在一个邻域 \(V\)（\(e\) 的），使得 \(F(x) > 0\)（对 \(x \in V\)）。现在，

\[
\mathrm{F} (x) = \int \varphi_ {\mathrm{B}} (s) \varphi_ {\mathrm{B}} (x ^ {- 1} s) d \beta (s) = \beta (\mathrm{B} \cap x \mathrm{B}).
\]

因此，对 \(x \in V\)，有 \(B \cap xB \neq \varnothing\)，从而 \(x \in B \cdot B^{-1}\)。于是 \(V \subset B \cdot B^{-1} \subset A \cdot A^{-1}\)。

推论 1. --- 设 H 为 G 的一个子群，它对一个 Haar 测度 \(\beta\) 可测。则 H 或者是开的，或者是局部 \(\beta\)-可忽略的。

因为 \(H = H \cdot H^{-1}\)，所以若 H 不是局部 \(\beta\)-可忽略的，则 H 包含 e 的一个邻域（命题 17），从而是开的（GT, III, §2, No.~1 命题 4 的推论）。

推论 2. --- 设 L 为 G 的一个对乘法稳定的子集，其余集对一个 Haar 测度 \(\beta\) 局部可忽略。则 \(L = G\)。

因为 \(L^{-1}\) 与 \(L \cap L^{-1}\) 的余集都局部 \(\beta\)-可忽略。而 \(L \cap L^{-1}\) 是一个子群，故是开的（推论 1），从而是闭的。因此 \(G - (L \cap L^{-1})\)——它是开的且局部 \(\beta\)-可忽略的——是空的。于是 \(G = L \cap L^{-1}\)。

命题 18. --- 设 G 为局部紧群，\(\Gamma\) 为一个集合，其上带有一个乘法 \((u,v)\mapsto uv\) 与一个 Hausdorff 拓扑，使得：

\begin{enumerate}
\def\labelenumi{\arabic{enumi})}
\item
  \(\Gamma\) 的拓扑在平移下不变；
\item
  乘法在 \(\Gamma \times \Gamma\) 的每一紧子集上的限制是连续的。
\end{enumerate}

设 \(f: \mathbf{G} \to \Gamma\) 为 \(\mathbf{G}\) 到 \(\Gamma\) 中的一个映射，使得 \(f(xy) = f(x)f(y)\)（对 \(x, y\) 属于 \(\mathbf{G}\)），且对一个 Haar 测度 \(\beta\)（\(\mathbf{G}\) 上的）可测。则 \(f\) 连续。

置 \(g(x) = f(x^{-1})\)（对 \(x \in \mathbf{G}\)）。由于 \(f\) 与 \(g\) 都是 \(\beta\)-可测的，存在一个非 \(\beta\)-可忽略的紧子集 \(\mathbf{K}\)（\(\mathbf{G}\) 的），使得 \(f\) 与 \(g\) 在 \(\mathbf{K}\) 上的限制连续。映射 \((x,y) \mapsto f(xy^{-1}) = f(x)g(y)\)（从 \(\mathbf{K} \times \mathbf{K}\) 到 \(\Gamma\) 中）连续，因为 \(\Gamma\) 的乘法在 \(f(\mathbf{K}) \times g(\mathbf{K})\) 上连续；而这个映射可写成 \(\varphi \circ \psi\)，其中 \(\psi\) 是映射 \((x,y) \mapsto xy^{-1}\)（从 \(\mathbf{K} \times \mathbf{K}\) 到 \(\mathbf{K} \cdot \mathbf{K}^{-1}\) 上），\(\varphi\) 是 \(f\) 在 \(\mathbf{K} \cdot \mathbf{K}^{-1}\) 上的限制。设 \(\mathbf{R}\) 为在 \(\mathbf{K} \times \mathbf{K}\) 上由 \(\psi\) 定义的等价关系。映射 \(\psi'\)（\((\mathbf{K} \times \mathbf{K}) / \mathbf{R}\) 到 \(\mathbf{K} \cdot \mathbf{K}^{-1}\) 上的、由 \(\psi\) 过渡到商而导出的）连续，因此 \((\mathbf{K} \times \mathbf{K}) / \mathbf{R}\) 是 Hausdorff 的且 \(\psi'\) 是一个同胚。由于 \(\varphi \circ \psi\) 连续，可见 \(f\) 在 \(\mathbf{K} \cdot \mathbf{K}^{-1}\) 上的限制连续。而 \(\mathbf{K} \cdot \mathbf{K}^{-1}\) 是 \(e\) 的一个邻域（命题 17），因此 \(f\) 在 \(e\) 处连续。对每一 \(x_0 \in \mathbf{G}\)，\(f(x_0x) = f(x_0)f(x)\)，故 \(f\) 在 \(x_0\) 处连续，因为 \(\Gamma\) 的拓扑在平移下不变。

推论 1. --- 设 G 为局部紧群，\(\beta\) 为 G 上的一个 Haar 测度，E 为一个 Hausdorff 桶型局部凸空间，\(U\) 为 G 在 E 中的一个线性表示，使得 \(U(s) \in \mathcal{L}(\mathrm{E};\mathrm{E})\) 对所有 \(s \in \mathrm{G}\) 成立，且是 \(\beta\)-可测的（当 \(\mathcal{L}(\mathrm{E};\mathrm{E})\) 赋有点态收敛拓扑时）。则 \(U\) 是一个连续线性表示。

设 \(\Gamma\) 为 E 的自同构群，赋有点态收敛拓扑。这一拓扑是 Hausdorff 的且在平移下不变。设 K 为 \(\Gamma\) 的一个紧子集。则 K 在赋有点态收敛拓扑的 \(\mathcal{L}(\mathrm{E};\mathrm{E})\) 中有界，故是等度连续的（TVS, III, §4, No.~2 定理 1）；因此映射 \((u,v)\mapsto v\circ u\)（从 \(\mathbf{K}\times \mathbf{K}\) 到 \(\mathcal{L}(\mathbf{E};\mathbf{E})\) 中）连续（同处, §5, No.~5 命题 9 的推论 1）。于是，对每一 \(x\in \mathbf{E}\)，G 到 E 中的映射 \(s\mapsto U(s)x\) 连续（命题 18）。由于 E 是桶型的，\(U\) 连续（§2, No.~1 命题 1）。

推论 2. --- 设 G 为局部紧群，\(\beta\) 为 G 上的一个 Haar 测度，E 为一个可分 Banach 空间，\(U\) 为 G 在 E 中的一个线性表示，使得 \(U(s) \in \mathcal{L}(\mathrm{E};\mathrm{E})\) 对所有 \(s \in \mathrm{G}\) 成立。设 \((a_m)\) 为 E 中的一个完全序列，\((a_n')\) 为 E 的对偶 E' 的单位球 B'（赋弱拓扑）中的一个稠密序列。假设 G 上的函数 \(s \mapsto \langle U(s)a_m, a_n' \rangle\) 都是 \(\beta\)-可测的。则 \(U\) 是一个连续线性表示。

我们先证明，对每一 \(z' \in E'\)，纯量函数

\[
s \mapsto \langle U (s) a _ {m}, z ^ {\prime} \rangle
\]

是 \(\beta\)-可测的；我们可以限制在 \(\|z'\| \leqslant 1\) 的情形，并且由于 \(B'\) 对弱拓扑是可度量的（TVS, III, §3, No.~4 命题 6 的推论 2），存在一个子列 \((a_{n_k}')\)（\((a_n')\) 的），它弱收敛于 \(z'\)；函数

\[
s \mapsto \langle U (s) a _ {m}, z ^ {\prime} \rangle
\]

从而是一列 \(\beta\)-可测函数的极限，由此得出我们的断言。由此推出，映射 \(s \mapsto U(s)a_{m}\)（G 到 E 中的）对每个 m 是 \(\beta\)-可测的（Ch. IV, §5, No.~5 命题 10）。另一方面，存在 E 的元素的一个序列 \((b_{m})\)——它们是 \(a_{i}\) 的线性组合——在 E 的单位球中稠密。对每一 \(s \in \mathbf{G}\)，\(\|U(s)\| = \sup_{m} \|U(s)b_{m}\|\)，因此 \(s \mapsto \|U(s)\|\) 可测。设 K 为 G 的一个紧子集，且设 \(\varepsilon > 0\)。存在 K 的一个紧子集 \(\mathrm{K}_{0}\)，使得 \(\beta(\mathrm{K} - \mathrm{K}_{0}) \leqslant \varepsilon\)，且函数在 \(\mathrm{K}_{0}\) 上的限制连续——这里函数指 \(s \mapsto U(s)a_{m}\) 与 \(s \mapsto \|U(s)\|\)。于是诸 \(U(s)\)（对 \(s \in \mathrm{K}_{0}\)）是等度连续的，且点态收敛拓扑在 \(U(\mathrm{K}_{0})\) 上诱导出在 \(a_{m}\) 组成的集合上的点态收敛拓扑（TVS, III, §3, No.~4 命题 5）。因此映射 \(s \mapsto U(s)\)（从 \(\mathrm{K}_{0}\) 到 \(\mathcal{L}_{s}(\mathrm{E};\mathrm{E})\) 中）连续。然后只需应用推论 1。

\subsection{正则化}

命题 19. --- 设 G 为局部紧群，\(\beta\) 为 G 上一个 \(\neq 0\) 的相对不变正测度，B 为 \(e\) 在 G 中的邻域滤子的一个基，由紧邻域组成。对每个 V \(\in\) B，设 \(f_{V}\) 为 G 上的一个连续函数 \(\geqslant 0\)，其支集包含在 V 中，且满足 \(\int f_{V} d\beta = 1\)。若 \(\mu\) 是 G 上的一个测度，则在 \(\mathcal{M}(G)\)（赋有 \(\mathcal{K}(G)\) 中的紧收敛拓扑）中，

\[
\mu = \lim _ {\mathrm{V}} (\mu * f _ {\mathrm{V}}) \cdot \beta = \lim _ {\mathrm{V}} (f _ {\mathrm{V}} * \mu) \cdot \beta ,
\]

极限是沿 B 的截滤子取的。

对 \(\mathcal{C}(\mathrm{G})\) 中的紧收敛拓扑，\(f_{V} \cdot \beta\) 沿 B 的截滤子趋于 \(\varepsilon_{e}\)（\(\S2\), No.~7 引理 4 的推论 1）。因此 \(\mu = \lim_{\mathrm{V}} \mu * (f_{\mathrm{V}} \cdot \beta) = \lim_{\mathrm{V}} (f_{\mathrm{V}} \cdot \beta) * \mu\) 在 \(\mathcal{M}(\mathrm{G})\) 中成立，该空间赋有 \(\mathcal{K}(\mathrm{G})\) 中的紧收敛拓扑（\(\S3\), No.~3 命题 12 的推论）。

备注. --- 1) 由此我们看到，G 上的每个测度都是一些测度的极限，这些测度关于每个 Haar 测度都具有连续密度（对命题 19 中所指的拓扑，更不用说对淡拓扑了）。

\begin{enumerate}
\def\labelenumi{\arabic{enumi})}
\setcounter{enumi}{1}
\tightlist
\item
  若 G 可度量化，则 B 可取为一列邻域 \((\mathrm{V}_{n})\)。这时 \(\mu\) 是具有连续密度的测度序列 \((\mu * f_{\mathrm{V}_{n}}) \cdot \beta\) 的极限。*若 G 是一个实 Lie 群，则 \(f_{V_{n}}\) 可取为无穷次可微的；我们以后会看到，密度 \(\mu * f_{V_{n}}\) 这时是无穷次可微的。*
\end{enumerate}

命题 20. --- 我们保留命题 19 的假设与记号。设 \(p \in [1, +\infty[\) 与 \(g \in \mathrm{L}^p(\mathbf{G}, \beta)\)。则

\[
g = \lim _ {\mathrm{V}} g * ^ {\beta} f _ {\mathrm{V}} = \lim _ {\mathrm{V}} f _ {\mathrm{V}} * ^ {\beta} g
\]

在范数 \(N_{p}\) 的意义下，极限是沿 B 的截滤子取的。

只需应用命题 6 (iii) 以及 §2, No.~7 引理 4 的推论 3。

备注 3. --- 由命题 15，函数 \(g * f_{\mathrm{V}}\)、\(f_{\mathrm{V}} * g\) 属于 \(\overline{\mathcal{K}(\mathrm{G})}\)

推论. --- 设 W 为 \(\mathrm{L}^{1}(\mathrm{G},\beta)\) 的一个闭线性子空间。要使 W 是 \(\mathrm{L}^{1}(\mathrm{G},\beta)\) 的一个左（相应地，右）理想，必须且只需 W 在 G 的左（相应地，右）平移下不变。

假设 W 是一个左理想。设 \(s \in G\) 与 \(g \in W\)。我们有 \(\varepsilon_{s} * g = \lim_{\mathbb{V}} f_{\mathbb{V}} * (\varepsilon_{s} * g) = \lim_{\mathbb{V}} (f_{\mathbb{V}} * \varepsilon_{s}) * g\)，而 \((f_{\mathbb{V}} * \varepsilon_{s}) * g \in \mathbb{W}\)，因此 \(\varepsilon_{s} * g \in W\)，从而 \(\boldsymbol{\gamma}(s)g \in \mathbb{W}\)。反之，若 W 在左平移下不变，则 \(\mu *^{\beta}g \in W\)（对 \(\mu \in \mathcal{M}^{1}(\mathbf{G})\) 与 \(g \in W\)），因此 W 更是 \(\mathrm{L}^{1}(\mathrm{G}, \beta)\) 的一个左理想。对右理想可类似论证。

例. --- 我们取 \(\mathbf{G} = \mathbf{R}\)。我们用下式定义一个函数 \(\mathbf{F}_n \in \mathcal{K}(\mathbf{R})\)

\[
\begin{array}{l l} F _ {n} (x) = (1 - x ^ {2}) ^ {n} & \text { if } x \in [ - 1, 1 ] \\ F _ {n} (x) = 0 & \text { if } x \notin [ - 1, 1 ]. \end{array}
\]

设 \(\mathbf{A}_n = \int_{-1}^{+1}\mathbf{F}_n(x)dx\)，且 \(\mathbf{G}_n = \mathbf{A}_n^{-1}\mathbf{F}_n\)。立即可见测度 \(\mathbf{G}_n(x)dx\) 满足 §2, No.~7 引理 4 的推论 1 的条件。设 \(\mu\) 为 \(\mathbf{R}\) 上的一个测度，其支集包含在 \([-1/2, 1/2]\) 中。则

\[
\begin{array}{r l} & {(\mu * \mathrm{G} _ {n}) (x) = \int_ {\mathbf {R}} \mathrm{G} _ {n} (x - y) d \mu (y)} \\ & {\qquad = \mathrm{A} _ {n} ^ {- 1} \int_ {- 1 / 2} ^ {1 / 2} \mathrm{F} _ {n} (x - y) d \mu (y).} \end{array}
\]

若 \(-1 / 2\leqslant x\leqslant 1 / 2\)，则

\[
(\mu * \mathrm{G} _ {n}) (x) = \mathrm{A} _ {n} ^ {- 1} \int_ {- 1 / 2} ^ {1 / 2} [ 1 - (x - y) ^ {2} ] ^ {n} d \mu (y),
\]

因此 \(\mu * \mathbf{G}_n\) 在 \([-1/2, 1/2]\) 上与一个多项式重合。特别地，若 \(f\) 是一个连续函数，其支集包含在 \([-1/2, 1/2]\) 中，则 \(f * \mathbf{G}_n\) 在 \([-1/2, 1/2]\) 中与一个多项式重合；此外，由命题 5 (iv) 以及 §2, No.~7 引理 4 的推论 3，\(f * \mathbf{G}_n\) 一致收敛于 \(f\)。*若 \(f\) 是 \(\mathbf{C}^r\) 类的，则导数 \(\mathrm{D}^s(f * \mathbf{G}_n)\) 一致趋于 \(\mathrm{D}^s f\)（对 \(0 \leqslant s \leqslant r\)）。*

\section{闭子群的空间}

在本节中，G 表示一个局部紧群，\(\mu\) 表示 G 上的一个右 Haar 测度。

\subsection{G 的闭子群上的 Haar 测度所成的空间}

引理 1. --- 设 \(\alpha\) 为 G 上一个 \(\neq 0\) 的正测度，S 为它的支集；下列两个条件等价：

\begin{enumerate}
\def\labelenumi{\alph{enumi})}
\item
  S 是 G 的一个闭子群，且 \(\alpha\) 在 S 上诱导的测度是 S 上的一个右 Haar 测度。
\item
  \(\delta(s)\alpha = \alpha\) 对每一 \(s \in S\) 成立。
\end{enumerate}

此外，当这些条件满足时，\(t \in \mathbf{G}\) 中使得 \(\delta(t)\alpha = \alpha\) 的那些 t 组成的集合等于 \(S\)。

显然 a) 蕴含 b)；反之，关系 b) 蕴含 \(Sx = S\)（对每一 \(x \in S\)）；换句话说，关系 \(x \in S\) 与 \(y \in S\) 蕴含 \(y \in Sx\)，或者再说一次 \(yx^{-1} \in S\)，而由于 \(S\) 非空，\(S\) 是 \(G\) 的一个闭子群。于是 \(t \in G\) 中使得 \(St = S\) 的那些 t 组成的集合等于 \(S\) 自身，由此得出最后的断言。

在本节其余部分，我们用 \(\Gamma\) 表示 G 上满足引理 1 的条件的 \(\neq0\) 正测度组成的集合，且对每一 \(\alpha\in\Gamma\)，用 \(H_{\alpha}\) 表示作为 \(\alpha\) 的支集的 G 的闭子群。

命题 1. --- 集合 \(\Gamma\) 在赋有淡拓扑的空间 \(\mathcal{M}_{+}(G) = \{0\}\) 中是闭的。

我们先证明下列引理：

引理 2. --- 设 X 为一个局部紧空间，且对每一测度 \(\alpha \in \mathcal{M}_{+}(X) - \{0\}\)，设 \(S_{\alpha}\) 为 \(\alpha\) 的支集。设 \(\Phi\) 为 \(\mathcal{M}_{+}(X) - \{0\}\) 上的一个滤子，它淡收敛于一个测度 \(\alpha_{0} \neq 0\)。则对 \(\alpha_{0}\) 的支集的一点 s 的每个邻域 V，存在一个集合 \(M \in \Phi\)，使得对每一 \(\alpha \in M\)，有 \(V \cap S_{\alpha} \neq \varnothing\)。

因为，若 \(\varphi \in \mathcal{K}_{+}(\mathbf{X})\) 是一个支集包含在 \(\mathbf{V}\) 中且满足 \(\int \varphi(x)d\alpha_0(x) > 0\) 的函数，则按定义存在一个集合 \(\mathbf{M} \in \Phi\)，使得 \(\int \varphi(x)d\alpha(x) > 0\)（对所有 \(\alpha \in \mathbf{M}\)），这就蕴含 \(\mathbf{V} \cap \mathbf{S}_{\alpha} \neq \emptyset\)。

引理 3. --- 设 E 为一个由滤子 \(\Phi\) 加滤的集合，且设 \(\xi \mapsto \alpha(\xi)\) 为 E 到 \(\Gamma\) 中的一个映射，它关于 \(\Phi\) 淡收敛于一个测度 \(\alpha_0 \neq 0\)。另一方面，设 \(\xi \mapsto t_{\xi}\) 为 E 到 G 中的一个映射，使得 \(t_{\xi} \in \mathrm{H}_{\alpha(\xi)}\)（对每一 \(\xi \in \mathrm{E}\)）。若 \(s\) 是映射 \(\xi \mapsto t_{\xi}\) 关于 \(\Phi\) 的一个聚点，则 \(\delta(s)\alpha_0 = \alpha_0\)。

若有必要，把 \(\Phi\) 换成一个更细的滤子，我们可以假设 \(s\) 是 \(\xi \mapsto t_{\xi}\) 关于 \(\Phi\) 的一个极限；由引理 1，\(\delta(t_{\xi})\alpha(\xi) = \alpha(\xi)\)（对每一 \(\xi \in E\)），而结论由映射 \((u,\lambda) \mapsto \delta(u)\lambda\) 在 \(G \times \mathcal{M}_{+}(G)\) 上的连续性得出（§3, No.~3 命题 13）。

为证明命题 1，由引理 1，只需证明：若一个滤子 \(\Psi\)（\(\Gamma\) 上的）淡收敛于一个测度 \(\alpha_0 \neq 0\)，且 \(s\) 属于 \(\alpha_0\) 的支集，则 \(\delta(s)\alpha_0 = \alpha_0\)。现在，对每一邻域 \(V\)（\(s\) 在 \(G\) 中的），由引理 2，存在一个 \(M \in \Psi\)，使得对每一 \(\alpha \in M\)，有 \(V \cap H_\alpha \neq \varnothing\)。对每一邻域 \(V\)（\(s\) 的）与每一 \(\alpha \in \Gamma\)，设 \(t_{V,\alpha}\) 为 \(V \cap H_\alpha\) 的一点（若 \(V \cap H_\alpha \neq \varnothing\)），在相反情形设为 \(H_\alpha\) 的任意一点；若 \(\Theta\) 是 \(s\) 的邻域滤子的截滤子，且 \(\Phi\) 是积滤子 \(\Theta \times \Psi\)，则由上所述，\(s\) 是 \((V,\alpha) \mapsto t_{V,\alpha}\) 关于 \(\Phi\) 的一个聚点。由于另一方面映射 \((V,\alpha) \mapsto \alpha\) 以 \(\alpha_0\) 为极限（关于 \(\Phi\)），命题由引理 3 得出。

命题 2. --- 设 \(\varphi\) 为 \(\mathcal{K}_{+}(\mathbf{G})\) 中的一个函数，使得 \(\varphi(e) > 0\)。则集合 \(\Gamma_{\varphi}\)（即测度 \(\alpha \in \Gamma\) 中满足 \(\int \varphi(x) d\alpha(x) = 1\) 的那些）对淡拓扑是紧的。

集合 \(\Gamma_{\varphi}\) 是 \(\Gamma\) 与 \(\mathcal{M}(\mathrm{G})\) 中超平面（由 \(\alpha\)——满足 \(\int \varphi(x)d\alpha(x)=1\)——形成）的交；由于这个超平面在 \(\mathcal{M}(\mathrm{G})\) 中是淡闭的且不含 0，由命题 1 可知 \(\Gamma_{\varphi}\) 在 \(\mathcal{M}(\mathrm{G})\) 中是淡闭的。因此只需证明：对 G 的每一紧子集 K，有 \(\sup_{\alpha\in\Gamma_{\varphi}}\alpha(\mathrm{K})<+\infty\)（Ch. III, §1,

No.~9, 命题 15)。现在，设 U 为 e 在 G 中由不等式 \(\varphi(x) > \varphi(e)/2\) 定义的开邻域；由于 \(1 = \int \varphi(x) d\alpha(x) \geqslant \int_{\mathrm{U}} \varphi(x) d\alpha(x)\)（对 \(\alpha \in \Gamma_{\varphi}\)），可见在置 \(c = 2/\varphi(e)\) 后，有 \(\alpha(\mathrm{U}) \leqslant c\)（对每一 \(\alpha \in \Gamma_{\varphi}\)）。设 V 为 e 在 G 中的一个对称开邻域，使得 \(V^{2} \subset U\)；我们来证明 \(\alpha(Vx) \leqslant c\)（对每一 \(x \in G\) 与每一 \(\alpha \in \Gamma_{\varphi}\)）。事实上，若 Vx 不与支集 \(H_{\alpha}\)（\(\alpha\) 的）相交，这一关系是平凡的；反之，若存在一个 \(h \in Vx \cap H_{\alpha}\)，则对某个 \(v \in V\) 有 h = vx，从而

\[
\mathbf {V} x = \mathbf {V} v ^ {- 1} h \subset \mathbf {V} ^ {2} h \subset \mathrm{U} h,
\]

而由于 \(\delta(h)\alpha = \alpha\)，得出 \(\alpha(\mathrm{V}x) \leqslant \alpha(\mathrm{U}h) = \alpha(\mathrm{U}) \leqslant c\)。现在设 \((x_i)_{1 \leqslant i \leqslant n}\) 为 K 的点的一个序列，使得诸 \(Vx_i\) 构成 K 的一个覆盖；由上所述，\(\alpha(K) \leqslant \sum_{i=1}^{n} \alpha(Vx_i) \leqslant nc\)（对每一 \(\alpha \in \Gamma_\varphi\)）；证毕。

命题 3. --- 在命题 2 的假设下，映射 \(\alpha \mapsto \left( \langle \varphi, \alpha \rangle, \frac{\alpha}{\langle \varphi, \alpha \rangle} \right)\) 是 \(\Gamma\) 到积空间 \(\mathbf{R}_{+}^{*} \times \Gamma_{\varphi}\) 上的一个同胚。

由于映射 \(\alpha \mapsto \langle \varphi, \alpha \rangle\) 是淡连续的，只需注意到 \(\langle \varphi, \alpha \rangle \neq 0\) 对每一测度 \(\alpha \in \Gamma\) 成立，因为 \(e\) 属于支集 \(\mathbf{H}_{\alpha}\)（\(\alpha\) 的）且 \(\varphi(e) > 0\)。

\subsection{齐性空间体积的半连续性}

在本 No. 中，对每一测度 \(\alpha \in \Gamma\)，我们置

\[
\mathrm{Q} _ {\alpha} = \mathrm{G} / \mathrm{H} _ {\alpha},\tag{1}
\]

且用 \(\pi_{\alpha}\) 表示典范映射 \(\mathbf{G} \to \mathbf{Q}_{\alpha}\)。

设 \(\Gamma^0\) 为 \(\Gamma\) 中由这样的测度 \(\alpha\) 组成的子集：子群 \(\mathbf{H}_{\alpha}\)（\(\mathbf{G}\) 的）是单模的；\(\Gamma^0\) 的元素由如下事实刻画：\(\alpha(f) = \alpha(\check{f})\) 对每一函数 \(f \in \mathcal{K}(\mathbf{G})\) 成立（\(\mathcal{K}(\mathbf{H}_{\alpha})\) 的每个函数都可由 Urysohn 定理延拓为 \(\mathcal{K}(\mathbf{G})\) 的函数）；由此得出 \(\Gamma^0\) 是 \(\Gamma\) 的一个闭子集。回忆：对每一 \(\alpha \in \Gamma^0\)，商测度 \(\mu_{\alpha} = \mu / \alpha\)（\(\mathbf{Q}_{\alpha}\) 上的）有定义且在 \(\mathbf{G}\) 下相对不变（Ch. VII, §2, No.~6 定理 3）；再回忆：对每一函数 \(f \in \mathcal{K}(\mathbf{G})\)，

\[
\int_ {\mathbf {G}} f (x) d \mu (x) = \int_ {\mathbf {Q} _ {\alpha}} d \mu_ {\alpha} (\dot {x}) \int_ {\mathrm{H} _ {\alpha}} f (x s) d \alpha (s),\tag{2}
\]

其中 \(\dot{x} = \pi_{\alpha}(x)\) 是 \(x\in \mathbf{G}\) 在 \(\mathbf{Q}_{\alpha}\) 中的典范像。

命题 4. --- 设 \(\Gamma^0\) 为测度 \(\alpha \in \Gamma\)（使得 \(\mathrm{H}_{\alpha}\) 是单模的）组成的集合，且对每一 \(\alpha \in \Gamma^0\) 置 \(\mu_{\alpha} = \mu / \alpha\)；则映射 \(\alpha \mapsto \| \mu_{\alpha} \|\)（从 \(\Gamma^0\) 到 \(\overline{\mathbf{R}}\) 中）对淡拓扑是下半连续的。

对每一 \(\alpha \in \Gamma^0\) 与每一函数 \(f\in \mathcal{K}(\mathbf{G})\)，置

\[
f _ {\alpha} (\dot {x}) = \int_ {\mathrm{H} _ {\alpha}} f (x s) d \alpha (s) = (f * \alpha) (x),
\]

其中卷积是相对于右 Haar 测度 \(\mu\) 取的，并且用到了 \(\check{\alpha} = \alpha\) 这一事实（§4, No.~4 公式 (11)）。我们知道（Ch. VII, §2, No.~1 命题 2），映射 \(f \mapsto f_{\alpha}\)（从 \(\mathcal{H}_{+}(G)\) 到 \(\mathcal{H}_{+}(Q_{\alpha})\) 中）是满射；因此，由 (2)，

\[
\| \mu_ {\alpha} \| = \sup _ {f \in \mathcal {K} _ {+} (\mathrm{G}), f \neq 0} \mu_ {\alpha} (f _ {\alpha}) / \| f _ {\alpha} \| = \sup _ {f \in \mathcal {K} _ {+} (\mathrm{G}), f \neq 0} \mu (f) / \| f _ {\alpha} \|,
\]

其中已置

\[
\left\| f _ {\alpha} \right\| = \sup _ {\dot {x} \in \mathrm{Q} _ {\alpha}} | f _ {\alpha} (\dot {x}) | = \sup _ {x \in \mathrm{G}} | (f * \alpha) (x) |.\tag{3}
\]

为建立命题，只需证明：给定 \(f \in \mathcal{K}_{+}(G)\)，映射 \(\alpha \mapsto \|f_{\alpha}\|\) 是淡连续的。现在，设 K 为 f 的支集；函数 \(f * \alpha\) 的支集包含在 \(KH_{\alpha}\) 中且在 \(H_{\alpha}\) 下右不变；从而

\[
\left\| f _ {\alpha} \right\| = \sup _ {x \in K} | (f * \alpha) (x) |.
\]

因此结论由如下事实得出：映射 \(\alpha\mapsto f*\alpha\)（从赋有淡拓扑的 \(\mathcal{M}_{+}(\mathrm{G})\) 到赋有紧收敛拓扑的 \(\mathcal{C}(\mathrm{G})\) 中）是连续的（§4, No.~2 备注 1）。

回忆：若对一个测度 \(\alpha \in \Gamma^0\)，\(\| \mu_{\alpha}\|\) 有限，则 G 必定是单模的（Ch. VII, §2, No.~6 定理 3 的推论 3）。

命题 5. --- 设 g 为一个 \(\mu\)-可积的正数值函数，且设 \(\Gamma^{0}(g)\) 为测度 \(\alpha\in\Gamma^{0}\)（使得 \(\int^{*}g(xs)d\alpha(s)\geqslant1\) 对所有 \(x\in G\)）组成的集合。则映射 \(\alpha\mapsto\|\mu_{\alpha}\|\)（从 \(\Gamma^{0}(g)\) 到 \(\overline{\mathbf{R}}\) 中）是淡连续的。

对每一测度 \(\alpha \in \Gamma^0 (\mathbf{G})\)，回忆（Ch. VII, §2, No.~3 命题 5）函数

\[
g _ {\alpha} (\dot {x}) = \int_ {\mathrm{H} _ {\alpha}} g (x s) d \alpha (s)
\]

\(\mu_{\alpha}\)-几乎处处（在 \(Q_{\alpha}\) 上）有定义，是 \(\mu_{\alpha}\)-可积的，且

\[
\int_ {G} g (x) d \mu (x) = \int_ {Q _ {\alpha}} g _ {\alpha} (\dot {x}) d \mu_ {\alpha} (\dot {x}).\tag{4}
\]

鉴于命题 4，只需证明：在 \(\Gamma^0(g)\) 中，\(\alpha \mapsto \| \mu_\alpha \|\) 是上半连续的。固定一个测度 \(\alpha \in \Gamma^0(g)\)，且设 \(K\) 为 \(G\) 的一个紧子集。在 \(Q_\alpha\) 上存在一个具有紧支集的连续函数，取值于 {[}0,1{]}，在紧集 \(\pi_\alpha(K)\) 上等于 1；由于映射 \(f \mapsto f_\alpha\)（从 \(\mathcal{H}_+(\mathrm{G})\) 到 \(\mathcal{H}_+(\mathrm{Q}_\alpha)\) 中）是满射（Ch. VII, §2, No.~1 命题 2），可见存在一个函数 \(f \in \mathcal{H}_+(\mathrm{G})\)，使得

\[
(f * \alpha) (x) = \int_ {\mathrm{G}} f (x s)   d \alpha (s) \left\{ \begin{array}{l} \leqslant 1 \text {for all} x \in \mathrm{G} \\ = 1 \text {for all} x \in \mathrm{K}. \end{array} \right.
\]

由于 \(\beta \mapsto f * \beta\) 是一个连续映射（从赋有淡拓扑的 \(\mathcal{M}_{+}(\mathbf{G})\) 到赋有紧收敛拓扑的 \(\mathcal{C}(\mathbf{G})\) 中）（§4, No.~2 备注 1），可见对每一 \(\varepsilon > 0\)，\(\mathrm{U}_{\varepsilon}\)（由 \(\beta \in \Gamma^0 (\mathbf{G})\) 中使得

\[
f _ {\beta} (\dot {x}) = \int_ {\mathrm{G}} f (x s) d \beta (s) > 1 - \varepsilon \quad \text { for   all } x \in \mathrm{K}
\]

的测度组成）是 \(\alpha\) 在 \(\Gamma^0 (g)\) 中的一个开邻域；于是对每一 \(\beta \in \mathrm{U}_{\varepsilon}\)，由公式 (2)，我们有

\[
\left\| \mu_ {\alpha} \right\| \geqslant \int_ {G} f (x) d \mu (x) = \int_ {Q _ {\beta}} f _ {\beta} (\dot {x}) d \mu_ {\beta} (\dot {x}) \geqslant (1 - \varepsilon) \mu_ {\beta} \left(\pi_ {\beta} (K)\right).\tag{5}
\]

给定一个数 \(\varepsilon > 0\)，选取一个函数 \(h \in \mathcal{H}_{+}(\mathrm{G})\)，使得 \(\int_{\mathrm{G}} |g(x) - h(x)| d\mu(x) \leqslant \varepsilon\)，并在上文中取 \(\mathrm{K} = \operatorname{Supp}(h)\)。对每一 \(\beta \in \Gamma^0(g)\)，按假设 \(g_\beta(\dot{x}) \geqslant 1\)（对 \(\mu_\beta\)）在 \(\mathbf{Q}_\beta\) 中几乎处处成立，因此

\[
\mu_ {\beta} \big (\mathrm{Q} _ {\beta} - \pi_ {\beta} (\mathrm{K}) \big) \leqslant \int_ {\mathrm{Q} _ {\beta} - \pi_ {\beta} (\mathrm{K})} g _ {\beta} (\dot {x}) d \mu_ {\beta} (\dot {x}) = \int_ {\mathrm{G} - \mathrm{KH} _ {\beta}} g (x) d \mu (x)
\]

由 (4)；由于 \(h\) 在 \(\mathbf{K}\) 之外为零，更不用说在 \(\mathrm{KH}_{\beta}\) 之外为零，由此得出

\[
\begin{array}{c} \mu_ {\beta} \big (\mathrm{Q} _ {\beta} - \pi_ {\beta} (\mathrm{K}) \big) \leqslant \int_ {\mathrm{G} - \mathrm{KH} _ {\beta}} | g (x) - h (x) | d \mu (x) \\ \leqslant \int_ {\mathrm{G}} | g (x) - h (x) | d \mu (x) \leqslant \varepsilon ; \end{array}
\]

把这一结果与 (5) 结合起来，可见

\[
\| \mu_ {\beta} \| \leqslant \varepsilon + \| \mu_ {\alpha} \| / (1 - \varepsilon)
\]

当 \(\beta \in \mathrm{U}_{\varepsilon}\) 时，这就完成了证明。

推论 1. --- 设 K 为 G 的一个紧子集，V 为 e 在 G 中的一个对称紧邻域，c 为一个实数 \textgreater{} 0。映射 \(\alpha \mapsto \|\mu_{\alpha}\|\) 在 \(\alpha \in \Gamma^0\)（满足 \(G = KH_{\alpha}\) 且 \(\alpha(V) \geqslant c\)）组成的集合上的限制是淡连续的。

因为，设 \(g \in \mathcal{K}_{+}(\mathbf{G})\) 为一个函数，使得 \(g(x) \geqslant 1 / c\)（对 \(x \in \mathrm{KV}\)）。对每一 \(x \in \mathbf{K}\)，

\[
\int g (x s) d \alpha (s) \geqslant \int_ {\mathrm{V}} g (x s) d \alpha (s) \geqslant 1
\]

对满足陈述中条件的 \(\alpha\)；此外，由于 \(\pi_{\alpha}(\mathbf{K}) = \mathrm{Q}_{\alpha}\)，有 \(\alpha \in \Gamma^0(g)\)，由此得出推论。

推论 2. --- 设 A 为 G 的一个 \(\mu\)-可积子集。映射 \(\alpha \mapsto \|\mu_{\alpha}\|\) 在集合 \(N_A\)——G 的离散子群 H 的规范化 Haar 测度中使得 \(G = AH\) 的那些——上的限制是淡连续的。

对 \(a \in \mathbf{A}\) 与 \(\alpha \in \mathbf{N}_{\mathbf{A}}\)，

\[
\int \varphi_ {\mathrm{A}} (a s) d \alpha (s) \geqslant \varphi_ {\mathrm{A}} (a) = 1,
\]

而由于 \(\pi_{\alpha}(\mathbf{A}) = \mathbf{Q}_{\alpha}\)，有 \(\mathrm{N}_{\mathrm{A}} \subset \Gamma^{0}(\varphi_{\mathrm{A}})\)，因此推论由命题 5 得出。

\subsection{G 的闭子群的空间}

我们用 \(\Sigma\) 表示 G 的闭子群组成的集合；若把每一测度 \(\alpha\in\Gamma\) 对应到子群 \(H_{\alpha}\)（即 \(\alpha\) 的支集），就得到 \(\Gamma\) 到 \(\Sigma\) 中的一个映射（称为典范映射），它显然是满射，并使我们可以典范地把 \(\Sigma\) 等同于 \(\Gamma\) 中比率为 \textgreater0 的位似群的同轨集。集合 \(\Sigma\)（赋有 \(\Gamma\) 上淡拓扑的商拓扑）称为 G 的闭子群的空间。

定理 1. --- 设 G 为局部紧群。G 的闭子群组成的空间 \(\Sigma\) 是紧的。此外，我们有下列性质：

\begin{enumerate}
\def\labelenumi{(\roman{enumi})}
\item
  \(\Sigma^0\)（\(\mathbf{G}\) 的单模闭子群组成的集合）在 \(\Sigma\) 中是闭的（从而是紧的）。
\item
  若 G 由 e 的一个紧邻域生成，则集合 \(\Sigma_{c}^{0}\)——由 G 中使得商空间 G/H 为紧的单模闭子群 H 组成——在 \(\Sigma^{0}\) 中是开的（从而是局部紧的）。
\item
  对 e 在 G 中的每一相对紧开邻域 U，集合 \(D_{U}\)——由 G 中使得 \(H \cap U = \{e\}\) 的离散子群 H 组成——在 \(\Sigma^{0}\) 中是闭的（从而是紧的）。
\end{enumerate}

由 No.~1 的命题 3 可知 \(\Sigma\) 同胚于 \(\Gamma_{\varphi}\)，从而由 No.~1 的命题 2 是紧的。此外，在 No.~2 开头已指出，集合 \(\Gamma^0\)——由 \(\alpha \in \Gamma\) 中使得 \(\mathbf{H}_{\alpha}\) 单模的测度组成——在 \(\Gamma\) 中是闭的；由于 \(\Gamma^0\) 在比率为 \(>0\) 的位似下稳定，像 \(\Sigma^0\)（\(\Gamma^0\) 在 \(\Sigma\) 中的）是 \(\Sigma\) 的一个闭子集，这就证明了 (i)。

性质 (ii) 将是下列命题的一个推论：

命题 6. --- 假设局部紧群 G 由 e 的一个紧邻域生成。则集合 \(\Gamma_{c}^{0}\)——由测度 \(\alpha\in\Gamma^{0}\)（使得 \(\mathbf{G} / \mathbf{H}_{\alpha}\) 为紧的）组成——在 \(\Gamma^0\) 中是开的，且 \(\Gamma_c^0\) 上映射 \(\alpha \mapsto \| \mu_{\alpha}\|\) 的限制是淡连续的。

用 No.~2 命题 5 的记号，我们有，对 \(g \in \mathcal{K}_{+}(\mathbf{G})\)，

\[
\Gamma^ {0} (g) \subset \Gamma_ {c} ^ {0}.\tag{6}
\]

因为，若 K 为 g 的支集，则关系 \(\int g(xs) d\alpha(s) \geqslant 1\)（对所有 \(x \in G\)）蕴含 \(KH_{\alpha} = G\)，因为积分在 \(KH_{\alpha}\) 的余集上显然为零，于是 \(G/H_{\alpha} = \pi_{\alpha}(K)\) 是紧的。给定一个测度 \(\alpha \in \Gamma_{c}^{0}\)，因此只需定义一个函数 \(g \in \mathcal{K}_{+}(G)\)，使得 \(\Gamma^{0}(g)\) 是 \(\alpha\) 在 \(\Gamma^{0}\) 中的一个邻域。由于 \(G/H_{\alpha}\) 是紧的，且典范映射 \(f \mapsto f_{\alpha}\)（从 \(\mathcal{K}_{+}(G)\) 到 \(\mathcal{K}_{+}(G/H_{\alpha})\) 中）是满射（Ch. VII, §2, No.~2），存在一个函数 \(g \in \mathcal{K}_{+}(G)\)，使得 \(\int g(xs) d\alpha(s) = 2\) 对每一 \(x \in G\) 成立。设 K 为 g 的（紧）支集，L 为 e 在 G 中的一个生成 G 的对称紧邻域；映射 \(\beta \mapsto g*\beta\)（从 \(\mathcal{M}_{+}(G)\) 到 \(\mathcal{C}(G)\) 中）是淡连续的（§4, No.~2 备注 1），故存在 \(\alpha\) 在 \(\Gamma^{0}\) 中的一个邻域 W，使得

\[
(g * \beta) (x) = \int g (x s) d \beta (s) \geqslant 1\tag{7}
\]

对所有 \(\beta \in W\) 与 \(x \in LK\) 成立。(7) 的左端在 \(\mathrm{KH}_{\beta}\) 之外为零，故关系 \(\beta \in W\) 蕴含

\[
\mathbf {L K} \subset \mathbf {K H} _ {\beta},
\]

由此通过对 n 作归纳推出，\(L^{n}K \subset KH_{\beta}\) 对每个整数 n \textgreater{} 0 成立；由于 L 生成 G，我们有 \(G = KH_{\beta}\)（对每一测度 \(\beta \in W\)），这就证明了 \(W \subset \Gamma_{c}^{0}\)。另一方面，(7) 的左端在 \(H_{\beta}\) 下右不变，故不等式 (7) 对 \(x \in LKH_{\beta} = G\) 也成立；因此 \(W \subset \Gamma^{0}(g)\)，这就证明了命题。

最后，(iii) 将是下列命题的一个推论：

命题 7. --- 设 \(N \subset \Gamma^{0}\) 为 G 的离散子群上的规范化 Haar 测度组成的子空间，且对 e 在 G 中的每一相对紧开邻域 U，设 \(N_{U}\) 为 N 中由 \(\alpha\)（满足 \(H_{\alpha} \cap U = \{e\}\)）组成的子集。则：

\begin{enumerate}
\def\labelenumi{\alph{enumi})}
\item
  \(\mathbf{N}_{\mathrm{U}}\) 是紧的。
\item
  \(\mathbf{N}_{\mathbf{U}}\) 在 \(\mathbf{N}\) 中的内部构成 \(\mathbf{N}\) 的一个覆盖（当 \(\mathbf{U}\) 取遍 \(e\) 在 \(\mathbf{G}\) 中的相对紧开邻域组成的集合时）。
\item
  要使 N 的一个子集 M 在 N 中相对紧，必须且只需存在 e 在 G 中的一个相对紧开邻域 U，使得 \(M \subset N_{U}\)。
\end{enumerate}

由于 \(D_{U}\) 是 \(N_{U}\) 在典范连续映射 \(\Gamma \rightarrow \Sigma\) 下的像，定理 1 的断言 (iii) 立即由命题 7 a) 得出。

为证明命题 7，我们注意到 \(N_{U}\) 可定义为 \(\Gamma^{0}\) 中由满足下列二条件的 \(\alpha\) 组成的子集

\[
\alpha (\{e \}) \geqslant 1 \quad \text { and } \quad \alpha (U) \leqslant 1.
\]

现在，若 A 在 G 中是紧的（相应地，是开的且相对紧的），则映射 \(\alpha \mapsto \alpha(A)\)（从 \(\mathcal{M}_{+}(G)\) 到 R 中）对淡拓扑是上半（相应地，下半）连续的（Ch. IV, §4, No.~4 命题 5 的推论 3 以及同处, §1, No.~1 命题 4）；由此可见 \(N_{U}\) 是 \(\Gamma^{0}\) 的一个闭子集。此外，设 \(\varphi \in \mathcal{K}_{+}(G)\) 为一个函数，使得 \(\varphi(e) = 1\) 且在 G - U 上 \(\varphi(x) = 0\)；显然 \(\int \varphi(x)d\alpha(x) = 1\) 对所有 \(\alpha \in N_{U}\) 成立；于是 No.~1 的命题 2 表明 \(N_{U}\) 是一个紧集，这就证明了 a)。另一方面，设 V 为 \(e\) 在 G 中的一个相对紧开邻域，使得 \(\overline{V} \subset U\)，且设 \(\varphi \in \mathcal{K}_{+}(G)\)，其支集包含在 U 中且在 V 上 \(\varphi(x) = 1\)。则 \(\alpha(\varphi) = 1\)（对 \(\alpha \in N_{U}\)），因此存在 \(\alpha\) 在 N 中的一个邻域 W，使得 \(\beta(\varphi) < 2\)（对 \(\beta \in W\)）；于是显然 \(W \subset N_{V}\)，因此 \(N_{V}\) 是 \(N_{U}\) 的一个邻域。由于诸 \(N_{U}\) 覆盖 N，这就证明了 b)。最后，N 的每一紧子集 M 都包含在有限多个集合 \(N_{U_i}\)（\(1 \leqslant i \leqslant n\)）的并中，而由于 \(\bigcup_{i} N_{U_i} \subset N_{U}\)（其中 \(U = \bigcap_{i} U_i\)），这就证明了 c)。

推论. --- \(\Gamma^0\) 的子空间 N 是局部紧的。
